% =====================================================================
%  TP1 — Architecture Cloud Azure pour ShopEasy
%  Binôme : Louis SCARFONE & Maxence BOURRAGUE — Bloc 4 (Cloud Computing)
%  Moteur : LuaLaTeX (LuaHBTeX)
% =====================================================================
\documentclass[11pt,a4paper]{article}

% ---------- Langue & fontes ----------
\usepackage{fontspec}
\usepackage{polyglossia}
\setdefaultlanguage{french}
\setmainfont{TeX Gyre Pagella}
\setsansfont{TeX Gyre Heros}
\setmonofont{DejaVu Sans Mono}[Scale=0.80]
\usepackage{microtype}

% ---------- Respiration (look épuré) ----------
\usepackage{setspace}
\setstretch{1.08}
\setlength{\parindent}{0pt}
\setlength{\parskip}{0.7em plus 0.12em minus 0.05em}

% ---------- Géométrie ----------
\usepackage[a4paper,margin=2.2cm,headheight=16pt,headsep=14pt,footskip=20pt]{geometry}

% ---------- Couleurs (palette Azure) ----------
\usepackage{xcolor}
\definecolor{azure}{HTML}{0078D4}
\definecolor{azuredark}{HTML}{14375E}
\definecolor{azuremid}{HTML}{2E75B6}
\definecolor{azurelight}{HTML}{E8F1FB}
\definecolor{azurepale}{HTML}{F3F8FD}
\definecolor{okgreen}{HTML}{107C10}
\definecolor{okgreenbg}{HTML}{E8F5E9}
\definecolor{warnorange}{HTML}{B5520A}
\definecolor{warnbg}{HTML}{FBEEDF}
\definecolor{dangerred}{HTML}{A4262C}
\definecolor{codebg}{HTML}{F6F8FA}
\definecolor{codeframe}{HTML}{D8E1EC}
\definecolor{ink}{HTML}{1B2733}
\definecolor{rowalt}{HTML}{EBF3FC}
\color{ink}

% ---------- Divers ----------
\usepackage{graphicx}
\graphicspath{{../livrables/assets/}}
\usepackage{fontawesome5}
\usepackage{enumitem}
\setlist{leftmargin=1.4em,topsep=2pt,itemsep=2pt,parsep=0pt}
\usepackage{needspace}
\usepackage{fvextra}

% ---------- Emojis (Twemoji couleur) ----------
\usepackage{twemojis}
\newcommand{\emo}[1]{\raisebox{-0.14em}{\twemoji{#1}}}
\newcommand{\eTarget}{\emo{1f3af}}
\newcommand{\ePackage}{\emo{1f4e6}}
\newcommand{\eCheck}{\emo{2705}}
\newcommand{\eWarn}{\emo{26a0}}
\newcommand{\eBulb}{\emo{1f4a1}}
\newcommand{\eLock}{\emo{1f512}}
\newcommand{\eKey}{\emo{1f510}}
\newcommand{\eRocket}{\emo{1f680}}
\newcommand{\eChart}{\emo{1f4ca}}
\newcommand{\eMoney}{\emo{1f4b0}}
\newcommand{\eGear}{\emo{2699}}
\newcommand{\eGlobe}{\emo{1f310}}
\newcommand{\eCamera}{\emo{1f4f8}}
\newcommand{\ePin}{\emo{1f4cc}}
\newcommand{\eMag}{\emo{1f50d}}
\newcommand{\eFlag}{\emo{1f3c1}}
\newcommand{\eParty}{\emo{1f389}}
\newcommand{\eShield}{\emo{1f6e1}}
\newcommand{\eBroom}{\emo{1f9f9}}
\newcommand{\eFolder}{\emo{1f4c2}}
\newcommand{\eServer}{\emo{1f5a5}}
\newcommand{\eDB}{\emo{1f5c4}}
\newcommand{\eScale}{\emo{2696}}
\newcommand{\eRed}{\emo{1f534}}
\newcommand{\eOrange}{\emo{1f7e0}}
\newcommand{\eYellow}{\emo{1f7e1}}
\newcommand{\eGreen}{\emo{1f7e2}}
\newcommand{\eWorld}{\emo{1f30d}}
\newcommand{\arr}{\ensuremath{\rightarrow}}
\newcommand{\code}[1]{\texttt{#1}}

% ---------- Tableaux (tabularray) ----------
\usepackage{tabularray}
\UseTblrLibrary{booktabs}
\NewDocumentEnvironment{ltable}{m}{%
  \par\addvspace{4pt}\begin{tblr}{
    width=\linewidth,
    colspec={#1},
    row{1}={bg=azuredark,fg=white,font=\sffamily\bfseries\small},
    row{2-Z}={font=\small},
    row{even}={bg=rowalt},
    hline{1,Z}={1.1pt,solid,azuredark},
    hline{2}={0.5pt,solid,azuredark!45},
    colsep=7pt, rowsep=4pt,
    stretch=1.18,
  }%
}{%
  \end{tblr}\par\addvspace{4pt}%
}

% ---------- Boîtes (tcolorbox) ----------
\usepackage[most]{tcolorbox}
\tcbset{
  boxbase/.style={
    enhanced, breakable, sharp corners, boxrule=0pt, leftrule=4pt,
    left=4.5mm, right=4.5mm, top=3.6mm, bottom=3.6mm,
    before skip=12pt, after skip=12pt,
    fonttitle=\sffamily\bfseries, coltitle=white,
    attach boxed title to top left={xshift=0mm,yshift=-2.6mm},
    boxed title style={sharp corners, boxrule=0pt, left=2.5mm,right=2.5mm,top=0.9mm,bottom=0.9mm},
  },
}
\newtcolorbox{fiche}{boxbase, colback=azurelight, colframe=azure}
\newtcolorbox{notebox}[1][Point de vigilance]{boxbase, colback=warnbg, colframe=warnorange,
  colbacktitle=warnorange, title={\eWarn~#1}}
\newtcolorbox{tipbox}[1][Astuce]{boxbase, colback=azurepale, colframe=azuremid,
  colbacktitle=azuremid, title={\eBulb~#1}}
\newtcolorbox{etatbox}[1][État après l'atelier]{boxbase, colback=okgreenbg, colframe=okgreen,
  colbacktitle=okgreen, title={\eCheck~#1}}

% Bloc code / terminal
\fvset{breaklines=true, breakanywhere=true, fontsize=\footnotesize,
  breaksymbolleft=\raisebox{0.1ex}{\scriptsize\textcolor{azure}{$\hookrightarrow$}}}
\newtcolorbox{codebox}[1]{
  enhanced, breakable, sharp corners, boxrule=0.7pt,
  colback=codebg, colframe=codeframe,
  coltitle=azuredark, colbacktitle=codeframe, fonttitle=\ttfamily\bfseries\scriptsize,
  title={\faTerminal~~#1}, left=3mm, right=3mm, top=2.6mm, bottom=2.6mm,
  before skip=12pt, after skip=14pt,
}

% Capture d'écran encadrée + légende
\newcommand{\screenshot}[2]{%
  \par\addvspace{10pt}
  \begin{center}
  \begin{tcolorbox}[enhanced, sharp corners, boxrule=0.7pt, colframe=codeframe, colback=white,
      left=0.8mm, right=0.8mm, top=0.8mm, bottom=0.8mm, drop fuzzy shadow=black!18]
    \includegraphics[width=\linewidth]{#1}
  \end{tcolorbox}\par\vspace{4pt}
  {\small\itshape\color{black!58}#2}
  \end{center}\par\addvspace{12pt}
}

% Carte « chiffre-clé » pour la page de garde
\newcommand{\statcard}[2]{%
  \begin{tcolorbox}[enhanced, sharp corners, boxrule=0pt, leftrule=3pt, colframe=azure,
     colback=azurepale, width=0.305\linewidth, halign=center, valign=center, height=2.3cm,
     left=1mm, right=1mm, top=1mm, bottom=1mm, nobeforeafter]
     {\sffamily\bfseries\fontsize{22}{24}\selectfont\color{azuredark}#1}\\[1.5mm]
     {\sffamily\small\color{black!58}#2}
  \end{tcolorbox}%
}

% ---------- Sections (titlesec) ----------
\usepackage{titlesec}
\titleformat{\section}
  {\sffamily\LARGE\bfseries\color{azuredark}}
  {\tcbox[on line, sharp corners, boxrule=0pt, colback=azure, colupper=white,
     left=2.2mm,right=2.2mm,top=0.6mm,bottom=0.6mm]{\thesection}}
  {0.7em}{}[\vspace{1mm}{\color{azure}\titlerule[1.6pt]}]
\titlespacing*{\section}{0pt}{3.6ex plus 0.8ex}{2.0ex}
\titleformat{\subsection}
  {\sffamily\large\bfseries\color{azuremid}}{\thesubsection}{0.6em}{}
\titlespacing*{\subsection}{0pt}{2.8ex plus 0.5ex}{1.1ex}
\titleformat{\subsubsection}
  {\sffamily\bfseries\color{azuredark}}{}{0pt}{}
\titlespacing*{\subsubsection}{0pt}{1.9ex plus 0.3ex}{0.5ex}

% ---------- En-têtes / pieds (fancyhdr) ----------
\usepackage{fancyhdr}
\pagestyle{fancy}
\fancyhf{}
\fancyhead[L]{\sffamily\small\color{azuredark}\textbf{TP1 Azure}~\textcolor{azure}{\textbullet}~ShopEasy}
\fancyhead[R]{\sffamily\small\color{black!55}\nouppercase{\leftmark}}
\fancyfoot[C]{\sffamily\small\color{black!55}\thepage}
\fancyfoot[R]{\sffamily\scriptsize\color{black!42}Bloc 4 · Cloud Computing}
\fancyfoot[L]{\sffamily\scriptsize\color{black!42}SCARFONE · BOURRAGUE}
\renewcommand{\headrulewidth}{0pt}
\renewcommand{\footrulewidth}{0pt}
\renewcommand{\headrule}{\vspace{-2pt}{\color{azure!70}\rule{\headwidth}{1.1pt}}}

% ---------- Liens ----------
\usepackage[hidelinks]{hyperref}
\hypersetup{colorlinks=true, linkcolor=azuredark, urlcolor=azuremid,
  pdftitle={TP1 Azure - ShopEasy}, pdfauthor={Louis SCARFONE, Maxence BOURRAGUE}}

% Filet décoratif
\newcommand{\thinrule}{\par\vspace{1mm}{\color{azure!30}\hrule height 0.5pt}\vspace{1mm}}

\begin{document}

% =====================================================================
%  PAGE DE GARDE
% =====================================================================
\begin{titlepage}
\thispagestyle{empty}
\setlength{\parskip}{0pt}
\raggedright

\begin{tcolorbox}[enhanced, sharp corners, boxrule=0pt, colback=azuredark,
   width=\linewidth, left=13mm, right=13mm, top=14mm, bottom=14mm,
   borderline north={5pt}{0pt}{azure}, borderline south={2pt}{0pt}{azure}]
  \raggedright
  {\sffamily\bfseries\large\color{azure!75}\faMicrosoft~~MICROSOFT AZURE \quad\textbullet\quad BLOC 4 — CLOUD COMPUTING}

  \vspace{9mm}
  {\sffamily\bfseries\color{white}\fontsize{37}{42}\selectfont TP1 — Architecture Cloud Azure}

  \vspace{4mm}
  {\sffamily\color{azurelight}\LARGE pour l'application \textbf{ShopEasy}}

  \vspace{8mm}
  {\color{azure}\rule{5.2cm}{2.5pt}}

  \vspace{5mm}
  {\sffamily\color{white}\large Conception, déploiement via Azure CLI et analyse d'une architecture cloud cible.}
\end{tcolorbox}

\vspace{1.8cm}

\begin{minipage}[t]{0.5\linewidth}
  {\sffamily\bfseries\large\color{azuredark}Binôme}\\[3.5mm]
  {\large Louis \textsc{Scarfone}}\\[2mm]
  {\large Maxence \textsc{Bourragué}}
\end{minipage}\hfill
\begin{minipage}[t]{0.46\linewidth}
  {\sffamily\bfseries\large\color{azuredark}Cadre}\\[3.5mm]
  Mastère Dev Manager Full Stack — RNCP niveau 7\\[2mm]
  Cas fil rouge : \textbf{ShopEasy}\\[2mm]
  Azure for Students — région \texttt{swedencentral}
\end{minipage}

\vspace{1.9cm}

\noindent
\statcard{15}{ateliers + quiz}\hfill
\statcard{9}{services Azure}\hfill
\statcard{$\approx$\,52\,\$}{coût mensuel estimé}

\vfill
{\sffamily\large\bfseries\color{azuredark}Juin 2026}\hfill
{\sffamily\color{black!45}Mastère Dev Manager Full Stack}
\end{titlepage}

% =====================================================================
%  TABLE DES MATIÈRES
% =====================================================================
\thispagestyle{fancy}
{\sffamily\Huge\bfseries\color{azuredark}Sommaire}
\vspace{2mm}{\color{azure}\titlerule[1.6pt]}\vspace{4mm}
\hypersetup{linkcolor=azuredark}
{\setlength{\parskip}{1pt}\tableofcontents}
\newpage


% #####################################################################
%  ATELIER 1
% #####################################################################
\section{Atelier 1 — Analyse de l'existant}

\begin{fiche}
\textbf{\eTarget~Objectif} — identifier les limites de l'architecture actuelle et formuler les premiers besoins techniques.

\textbf{\ePackage~Livrable attendu} — un tableau d'analyse de l'existant et une liste courte des besoins techniques prioritaires.
\end{fiche}

\subsection{Rappel de l'architecture actuelle}
ShopEasy héberge une application interne de gestion de commandes (consultation clients, saisie de commandes,
génération de factures, dépôt de documents justificatifs) sur \textbf{un unique serveur physique}, dans les
locaux de l'entreprise :
\begin{itemize}
  \item Un serveur \textbf{Ubuntu unique} (mono-serveur).
  \item \textbf{Apache / Nginx} exposant l'application web.
  \item Application \textbf{PHP ou Node.js} installée \textbf{localement}.
  \item Base \textbf{MySQL installée sur le même serveur}.
  \item \textbf{Documents clients} stockés dans un \textbf{répertoire local}.
  \item \textbf{Sauvegardes manuelles et irrégulières}.
  \item \textbf{Aucun outil centralisé de supervision}.
  \item Un \textbf{compte administrateur partagé} par plusieurs personnes.
\end{itemize}

\subsection{Tableau d'analyse de l'existant}
\begin{tblr}{
  width=\linewidth, colspec={Q[l,0.6] X[l,1.2] X[l,1.5]},
  row{1}={bg=azuredark,fg=white,font=\sffamily\bfseries\small},
  column{1}={bg=azurelight,font=\bfseries},
  cell{1}{1}={bg=azuredark,fg=white},
  row{2-Z}={font=\small},
  hline{1,Z}={1.1pt,azuredark}, hline{2}={0.5pt,azuredark!45},
  hline{4,6,8,10,12}={0.5pt,azuredark!25},
  colsep=7pt, rowsep=4pt, stretch=1.2,
}
Domaine & Risque / limite identifiée & Impact possible pour l'entreprise \\
\SetCell[r=2]{m}Disponibilité
 & \textbf{Point de défaillance unique (SPOF)} : toute l'application repose sur un seul serveur.
 & Une panne matérielle/OS rend l'application \textbf{totalement indisponible} : plus de prise de commande, plus de facturation. \\
 & \textbf{Aucune redondance ni répartition de charge} (web + app + base colocalisés).
 & Aucune tolérance aux pannes ; un incident bloque simultanément les \textbf{trois couches} ; pics de charge non absorbés. \\
\SetCell[r=2]{m}Sécurité
 & \textbf{Compte administrateur partagé} entre plusieurs personnes.
 & \textbf{Perte de traçabilité} (impossible de savoir qui a agi), audit et conformité (RGPD) compromis, révocation difficile. \\
 & \textbf{Données + documents colocalisés} avec le serveur web, sans segmentation réseau ni chiffrement.
 & La \textbf{compromission du serveur web donne un accès direct} à la base MySQL et aux documents \arr{} \textbf{fuite de données}. \\
\SetCell[r=2]{m}Performance
 & \textbf{Ressources partagées} (CPU/RAM/disque/IO) entre web, applicatif et base sur une seule machine.
 & \textbf{Contention} des ressources, lenteurs sous charge, dégradation du temps de réponse. \\
 & \textbf{Pas de scalabilité} : verticale limitée, \textbf{horizontale impossible}.
 & Impossible d'absorber la \textbf{croissance} ou les \textbf{pics} (campagnes) \arr{} insatisfaction utilisateurs. \\
\SetCell[r=2]{m}Exploitation
 & \textbf{Aucune supervision centralisée} (ni métriques, ni logs, ni alertes).
 & Incidents \textbf{détectés tardivement}, diagnostic à l'aveugle, \textbf{MTTR élevé}. \\
 & \textbf{Administration 100\,\% manuelle}, non automatisée, non reproductible.
 & Opérations \textbf{sujettes aux erreurs humaines}, forte \textbf{dépendance aux personnes}. \\
\SetCell[r=2]{m}Coût
 & \textbf{Modèle CAPEX} : investissement matériel initial + \textbf{surdimensionnement} pour les pics.
 & Capital immobilisé, \textbf{capacité payée mais inutilisée} la majeure partie du temps. \\
 & \textbf{Coûts d'exploitation cachés} (maintenance, matériel, temps humain) \textbf{non mesurés}.
 & \textbf{Aucune visibilité financière} ni capacité de pilotage / d'optimisation. \\
\SetCell[r=2]{m}Sauvegarde
 & \textbf{Sauvegardes manuelles et irrégulières}.
 & \textbf{RPO non maîtrisé} : perte de données récentes (commandes, factures) probable. \\
 & \textbf{Pas de stratégie de restauration testée}, sauvegardes vraisemblablement \textbf{locales}.
 & \textbf{RTO inconnu} ; un sinistre du serveur/site entraîne une \textbf{perte totale} (pas de copie hors-site). \\
\end{tblr}

\subsection{Réponses aux questions guidées}

\textbf{1. Que se passe-t-il si le serveur physique tombe en panne ?}\\
Indisponibilité \textbf{totale et immédiate} : l'application web, la base et les documents étant sur la même
machine, aucune bascule n'est possible. L'activité commerciale s'arrête (plus de commandes, de factures, ni
d'accès aux dossiers clients). La reprise dépend du délai de réparation matériel et de la qualité des dernières
sauvegardes — donc d'un \textbf{RPO et d'un RTO aujourd'hui non maîtrisés}.

\textbf{2. Pourquoi la base de données locale représente-t-elle un risque ?}
\begin{itemize}
  \item \textbf{Couplage fort} : base et application se partagent CPU/RAM/disque \arr{} contention et SPOF commun.
  \item \textbf{Sécurité} : la compromission du serveur web donne un accès direct aux données.
  \item \textbf{Exploitation} : sauvegardes, haute disponibilité, patchs du moteur SQL à la charge de l'entreprise.
  \item \textbf{Pas d'isolation réseau} : la base n'est pas dans une couche dédiée et protégée.
\end{itemize}

\textbf{3. Quels problèmes pose l'absence de supervision ?}
\begin{itemize}
  \item Détection \textbf{tardive ou nulle} des incidents (on apprend la panne par les utilisateurs).
  \item Pas de métriques \arr{} \textbf{diagnostic à l'aveugle}, MTTR élevé.
  \item Impossible d'\textbf{anticiper la saturation} (disque plein, CPU, mémoire).
  \item Aucune alerte de sécurité (tentatives d'intrusion) ni de coût.
\end{itemize}

\textbf{4. Quels risques sont liés à un compte administrateur partagé ?}
\begin{itemize}
  \item \textbf{Perte de traçabilité / d'imputabilité} \arr{} audit et conformité compromis.
  \item \textbf{Non-respect du moindre privilège} : tout le monde dispose de tous les droits.
  \item \textbf{Difficulté de révocation} : un départ impose de changer le secret pour tous.
  \item \textbf{Surface d'attaque accrue} : secret partagé, pas de MFA individualisé.
\end{itemize}

\textbf{5. Quels éléments devraient être séparés dans une architecture cloud cible ?}
\begin{itemize}
  \item Les \textbf{couches} : réseau / applicatif (web) / données / stockage documentaire.
  \item La \textbf{segmentation réseau} : subnets distincts (web, data, admin) + filtrage \textbf{NSG}.
  \item La \textbf{base} \arr{} service managé isolé (Azure SQL Database), non exposé à Internet.
  \item Les \textbf{documents} \arr{} stockage objet dédié (Storage Account / Blob).
  \item Les \textbf{identités} \arr{} comptes nominatifs + \textbf{RBAC} (fin du compte partagé).
  \item La \textbf{disponibilité} \arr{} plusieurs instances derrière un \textbf{répartiteur de charge}.
  \item La \textbf{supervision} \arr{} plateforme dédiée (Azure Monitor).
\end{itemize}

\subsection{Besoins techniques prioritaires (synthèse)}
\begin{enumerate}
  \item \textbf{Supprimer le SPOF} : redondance applicative ($\geq 2$ instances) derrière un \textbf{répartiteur de charge}.
  \item \textbf{Séparer les couches} : segmentation réseau (\textbf{VNet + subnets} web/data) et filtrage (\textbf{NSG}).
  \item \textbf{Externaliser et fiabiliser la base} : base managée (Azure SQL Database), \textbf{non exposée}.
  \item \textbf{Externaliser les documents} : stockage objet durable (Storage Account / Blob) avec versioning.
  \item \textbf{Fiabiliser les sauvegardes} : stratégie automatisée et testée (RPO/RTO définis, copie hors-site).
  \item \textbf{Maîtriser les accès} : identités nominatives + \textbf{RBAC} (moindre privilège).
  \item \textbf{Superviser} : supervision centralisée avec \textbf{alertes} (Azure Monitor).
  \item \textbf{Piloter les coûts} : visibilité via \textbf{tags + Cost Management} (passage CAPEX \arr{} OPEX maîtrisé).
\end{enumerate}


% #####################################################################
%  ATELIER 2
% #####################################################################
\section{Atelier 2 — Choix des services Azure}

\begin{fiche}
\textbf{\eTarget~Objectif} — associer les besoins métier de ShopEasy aux services Azure pertinents.

\textbf{\ePackage~Livrable attendu} — un tableau besoins/services complété, avec justification technique et qualification IaaS / PaaS / gouvernance.
\end{fiche}

\subsection{Tableau de choix besoins \texorpdfstring{$\rightarrow$}{->} services}
\begin{ltable}{Q[l,0.8] Q[l,1.0] X[l,1.7] Q[c,0.7]}
Besoin & Service Azure choisi & Justification technique & Modèle \\
Hébergement de l'application & \textbf{Azure Virtual Machines} (2 VM Linux, Nginx) & Migration \textit{lift-and-shift} de l'app PHP/Node.js \textbf{sans refonte} ; contrôle de l'OS ; 2 instances pour la redondance. & \textbf{IaaS} \\
Répartition de charge & \textbf{Azure Load Balancer} (L4) & Distribue le trafic vers les 2 VM et retire une instance défaillante via sonde \arr{} supprime le SPOF. & IaaS / réseau \\
Base de données & \textbf{Azure SQL Database} & Base \textbf{managée} : sauvegardes, patchs, HA et chiffrement gérés par Azure ; remplace MySQL ; \textbf{non exposée}. & \textbf{PaaS} \\
Stockage des documents & \textbf{Azure Storage Account} (Blob) & Stockage objet \textbf{durable}, versioning et cycle de vie ; sort les fichiers du serveur ; accès \textbf{privé}. & \textbf{PaaS} \\
Isolation réseau & \textbf{Azure Virtual Network} (+ subnets) & Réseau privé \textbf{isolé} ; segmentation web / data / admin. & IaaS / réseau \\
Filtrage réseau & \textbf{Network Security Groups} & Filtrage L3/L4 par subnet ; n'ouvre que le nécessaire (80/443, 22 restreint, 1433 interne). & IaaS / réseau \\
Gestion des accès & \textbf{Microsoft Entra ID + RBAC} & Identités \textbf{nominatives} + MFA ; rôles au \textbf{moindre privilège} \arr{} fin du compte partagé. & Gouvernance \\
Monitoring & \textbf{Azure Monitor} & Métriques + logs, \textbf{alertes actionnables} \arr{} visibilité et détection précoce. & Gouvernance \\
Suivi des coûts & \textbf{Azure Cost Management} (+ tags) & Coûts réels, budgets et alertes, attribution par \textbf{tags} \arr{} pilotage FinOps. & Gouvernance \\
\end{ltable}

\subsection{Question de justification — IaaS / PaaS / gouvernance}

\subsubsection{Récapitulatif des qualifications}
\begin{ltable}{Q[l,0.45] X[l,1]}
Catégorie & Services concernés \\
\textbf{IaaS} (infrastructure) & Virtual Machines, Virtual Network, NSG, Load Balancer \\
\textbf{PaaS} (plateforme managée) & Azure SQL Database, Storage Account \\
\textbf{Gouvernance} (transverse) & Microsoft Entra ID + RBAC, Azure Monitor, Cost Management \\
\end{ltable}

\subsubsection{En quoi cela modifie la responsabilité de l'équipe}
Le choix IaaS / PaaS / gouvernance déplace le curseur du \textbf{modèle de responsabilité partagée} : plus un
service est managé, plus Azure prend en charge de couches techniques, et plus l'équipe se recentre sur la
\textbf{valeur métier}.
\begin{itemize}
  \item \textbf{IaaS} (VM, VNet, NSG, LB) — Azure gère le matériel ; \textbf{l'équipe reste responsable} de l'OS, des correctifs, du durcissement et du réseau. \arr{} contrôle maximal, charge d'exploitation la plus élevée.
  \item \textbf{PaaS} (SQL Database, Storage) — Azure gère en plus l'OS, le moteur, les patchs, la HA et les sauvegardes. \textbf{L'équipe se concentre} sur la donnée et la configuration. \arr{} administration fortement réduite.
  \item \textbf{Gouvernance} (Entra ID/RBAC, Monitor, Cost Management) — Azure fournit l'outil ; \textbf{l'équipe définit les politiques} (rôles, alertes, budgets). \arr{} responsabilité de pilotage et de conformité.
\end{itemize}

\subsubsection{Lecture pour ShopEasy}
En partant d'un existant \textbf{100\,\% à sa charge}, ShopEasy bascule vers un mix \textbf{IaaS + PaaS} qui
transfère à Azure l'administration de la base et du stockage, conserve la maîtrise de la couche applicative le
temps de la migration, et \textbf{ajoute une gouvernance} (identités, supervision, coûts) qui n'existait pas.

\begin{tipbox}[Suite logique]
Le TP2 réduira encore la part IaaS (App Service, Infrastructure as Code) une fois la migration stabilisée.
\end{tipbox}


% #####################################################################
%  ATELIER 3
% #####################################################################
\section{Atelier 3 — Architecture cible Azure}

\begin{fiche}
\textbf{\eTarget~Objectif} — construire une architecture Azure cible claire et justifiable pour ShopEasy.

\textbf{\ePackage~Livrable attendu} — un schéma d'architecture cible et une légende couvrant composants, flux, zones réseau, exposition Internet, composants internes, sécurité et supervision.
\end{fiche}

\subsection{Schéma d'architecture cible}
\screenshot{architecture.png}{Architecture cible : Internet \arr{} Load Balancer \arr{} 2 VM web (\code{snet-web}) ; VM \arr{} SQL Database et Storage en accès privé (\code{snet-data}) ; plan de gouvernance Entra ID/RBAC, Azure Monitor, Cost Management.}

\subsection{Légende — les 7 éléments demandés}

\subsubsection{1. Composants Azure}
\begin{ltable}{Q[l,0.62] X[l,1]}
Composant & Rôle \\
\textbf{Resource Group} \code{rg-shopeasy-dev} & Conteneur logique de toutes les ressources (cycle de vie, droits, coûts). \\
\textbf{Virtual Network} \code{vnet-shopeasy-dev} & Réseau privé isolé \code{10.10.0.0/16}. \\
\textbf{Subnets} web / data / admin & Séparation des couches applicative / données / administration. \\
\textbf{2 × Virtual Machines} & Couche applicative (Nginx + app), redondante. \\
\textbf{Azure Load Balancer} & Répartiteur de charge frontal avec sonde de santé. \\
\textbf{Azure SQL Database} & Base relationnelle managée (PaaS). \\
\textbf{Storage Account} & Stockage objet (Blob) des documents clients. \\
\textbf{Network Security Groups} & Filtrage réseau par subnet. \\
\textbf{Entra ID + RBAC} & Identités et droits d'administration. \\
\textbf{Azure Monitor} & Supervision (métriques, logs, alertes). \\
\textbf{Azure Cost Management} & Suivi et gouvernance des coûts. \\
\end{ltable}

\subsubsection{2. Flux principaux}
\begin{enumerate}
  \item \textbf{Internet \arr{} Load Balancer} : trafic web HTTP/HTTPS (80/443).
  \item \textbf{Load Balancer \arr{} VM 01 / 02} : répartition + sonde de santé (port 80).
  \item \textbf{VM \arr{} Azure SQL Database} : requêtes SQL (port \textbf{1433}, privé, depuis le subnet web).
  \item \textbf{VM \arr{} Storage Account} : lecture/écriture des documents (HTTPS, privé).
  \item \textbf{Admin \arr{} VM} : SSH (port 22) \textbf{restreint} (IP admin ou Azure Bastion).
  \item \textbf{Entra ID / RBAC \arr{} Resource Group} : contrôle des accès.
  \item \textbf{Toutes ressources \arr{} Azure Monitor} : remontée des métriques et journaux.
\end{enumerate}

\subsubsection{3. Zones réseau}
\begin{ltable}{Q[l,0.7] X[l,1] Q[l,0.7]}
Zone & Contenu & Exposition \\
\textbf{Frontal public} & Load Balancer (IP publique) & \textbf{Exposé Internet} (80/443) \\
\code{snet-web} & VM web 01 \& 02 & Interne (via le LB) \\
\code{snet-data} & Azure SQL Database, Storage (privé) & \textbf{Interne uniquement} \\
\code{snet-admin} & Azure Bastion (optionnel) & Administration contrôlée \\
\end{ltable}

\subsubsection{4. Composants exposés à Internet}
\begin{itemize}
  \item \textbf{Uniquement le Load Balancer} (80/443) — point d'entrée unique et contrôlé.
  \item \textit{En TP, les VM ont une IP publique temporaire pour le SSH restreint ; en cible production, on les en prive au profit d'Azure Bastion.}
\end{itemize}

\subsubsection{5. Composants internes (non exposés)}
\begin{itemize}
  \item VM web (joignables seulement via le Load Balancer).
  \item \textbf{Azure SQL Database} et \textbf{Storage Account} : accès \textbf{privé}, jamais directement depuis Internet.
\end{itemize}

\subsubsection{6. Éléments de sécurité}
\begin{itemize}
  \item \textbf{nsg-web} : autorise 80/443 depuis Internet, 22 depuis l'IP admin uniquement.
  \item \textbf{nsg-data} : autorise 1433 depuis \code{snet-web} uniquement \arr{} base non exposée.
  \item \textbf{Entra ID + RBAC} : comptes nominatifs, moindre privilège.
  \item SQL + Storage en accès privé (Private Endpoint en cible), chiffrement et TLS.
  \item \textbf{Segmentation réseau} en subnets pour limiter les mouvements latéraux.
\end{itemize}

\subsubsection{7. Éléments de supervision}
\begin{itemize}
  \item \textbf{Azure Monitor} : métriques VM (CPU/mémoire/disque), disponibilité du LB, métriques SQL et Storage, logs d'activité du Resource Group.
  \item \textbf{Alertes} : CPU élevé, échec des sondes, espace SQL proche de la limite, coût projeté supérieur au budget.
\end{itemize}

\subsection{Plan d'adressage IP}
\begin{ltable}{Q[l,0.5] Q[l,0.5] X[l,1]}
Élément & Plage CIDR & Rôle \\
\textbf{VNet} & \code{10.10.0.0/16} & Réseau global ShopEasy \\
\code{snet-web} & \code{10.10.1.0/24} & Machines virtuelles web \\
\code{snet-data} & \code{10.10.2.0/24} & Services de données / endpoints privés \\
\code{snet-admin} & \code{10.10.3.0/24} & Accès d'administration / bastion (optionnel) \\
\end{ltable}

\subsection{Variante « durcissement production »}
L'architecture ci-dessus est volontairement \textbf{lisible pour le TP}. Pour une mise en production réelle :
\begin{itemize}
  \item \textbf{Application Gateway + WAF} (L7) à la place du Load Balancer L4, avec terminaison TLS et HTTPS obligatoire.
  \item \textbf{Private Endpoint} pour Azure SQL et Storage (suppression de tout accès public).
  \item \textbf{Suppression des IP publiques des VM} + administration via Azure Bastion.
  \item \textbf{Déploiement multi-zones} (Availability Zones) des VM et du LB.
  \item \textbf{MFA} systématique et revue périodique des rôles RBAC.
  \item \textbf{Sauvegardes formalisées} + tests de restauration ; \textbf{Infrastructure as Code} (Terraform/Bicep) en TP2.
\end{itemize}


% #####################################################################
%  ATELIER 4
% #####################################################################
\section{Atelier 4 — Préparation de l'environnement Azure}

\begin{fiche}
\textbf{\eTarget~Objectif} — créer les conteneurs logiques et les conventions de nommage du projet.

\textbf{\ePackage~Livrable attendu} — capture du Resource Group et tableau de convention de nommage complété avec le suffixe.
\end{fiche}

\subsection{Convention de nommage adoptée}
Convention retenue : \code{<type>-<application>-<environnement>} (+ suffixe unique pour les ressources à nom
global). \textbf{Suffixe unique du projet : \code{ls01}} (initiales \textit{Louis Scarfone} + numéro).

\begin{ltable}{Q[l,0.7] Q[l,0.8] X[l,1.2]}
Ressource & Nom retenu & Remarque \\
Resource Group & \code{rg-shopeasy-dev} & — \\
Virtual Network & \code{vnet-shopeasy-dev} & — \\
Subnet web & \code{snet-web} & VM applicatives \\
Subnet data & \code{snet-data} & Services de données / endpoints privés \\
Subnet admin & \code{snet-admin} & Accès d'administration (optionnel) \\
NSG web & \code{nsg-web} & Filtrage du subnet web \\
NSG data & \code{nsg-data} & Filtrage du subnet data \\
Virtual Machine 1 & \code{vm-web-01} & — \\
Virtual Machine 2 & \code{vm-web-02} & — \\
\textbf{Storage Account} & \code{stshopeasyls01} & \textbf{Nom global unique}, minuscules, sans tiret \\
\textbf{Azure SQL Server} & \code{sql-shopeasy-ls01} & \textbf{Nom global unique} \\
Azure SQL Database & \code{sqldb-shopeasy} & — \\
\end{ltable}

\subsection{Création du Resource Group (Azure CLI)}
\begin{codebox}{bash}
\begin{Verbatim}
az group create \
  --name rg-shopeasy-dev \
  --location swedencentral \
  --tags projet=shopeasy environnement=dev module=cloud owner="Equipe Cloud" costcenter=Formation
\end{Verbatim}
\end{codebox}

\begin{notebox}[Contrainte régionale Azure for Students]
L'abonnement applique une policy \textit{« Allowed resource deployment regions »} qui restreint les
déploiements à 5 régions : \code{germanywestcentral}, \code{spaincentral}, \code{swedencentral},
\code{italynorth}, \code{uaenorth}. \textbf{\code{francecentral} est bloquée.} Après analyse (voir Atelier 7),
la seule région autorisée où les VM sont réellement déployables est \textbf{\code{swedencentral}} — c'est la
région retenue pour tout le TP, RG compris.
\end{notebox}

\subsubsection{Tags de gouvernance appliqués}
\begin{ltable}{Q[l,0.5] Q[l,0.6] X[l,1]}
Tag & Valeur & Usage \\
\code{projet} & \code{shopeasy} & Rattachement applicatif \\
\code{environnement} & \code{dev} & Distinction dev / prod \\
\code{module} & \code{cloud} & Contexte pédagogique \\
\code{owner} & \code{Equipe Cloud} & Responsable / contact \\
\code{costcenter} & \code{Formation} & Imputation des coûts (FinOps) \\
\end{ltable}

\subsection{Preuve de création (\code{az group show})}
\begin{codebox}{json}
\begin{Verbatim}
{
  "location": "swedencentral",
  "name": "rg-shopeasy-dev",
  "properties": { "provisioningState": "Succeeded" },
  "tags": {
    "costcenter": "Formation", "environnement": "dev", "module": "cloud",
    "owner": "Equipe Cloud", "projet": "shopeasy"
  },
  "type": "Microsoft.Resources/resourceGroups"
}
\end{Verbatim}
\end{codebox}

\screenshot{rg-global-view.png}{Vue d'ensemble du Resource Group \code{rg-shopeasy-dev} (région Sweden Central, abonnement Azure for Students).}
\screenshot{rg-tags.png}{Les 5 tags de gouvernance appliqués au Resource Group.}

\subsection{Questions de validation}

\textbf{1. Pourquoi regrouper les ressources dans un Resource Group dédié ?}\\
Un Resource Group matérialise un \textbf{cycle de vie commun} : toutes les ressources de ShopEasy se créent,
se gèrent, se sécurisent (RBAC au niveau du groupe), se suivent en coût et \textbf{se suppriment ensemble}.
Cela évite de mélanger des ressources de projets différents et simplifie l'exploitation, les droits et le
nettoyage de fin de TP (\code{az group delete}).

\textbf{2. Pourquoi ajouter des tags ?}
\begin{itemize}
  \item \textbf{Attribuer les coûts} (par application, environnement, centre de coût) dans Cost Management.
  \item \textbf{Rechercher et filtrer} les ressources.
  \item \textbf{Automatiser} des actions (extinction des ressources \code{environnement=dev}, Azure Policy).
  \item \textbf{Identifier le propriétaire} et la criticité.
\end{itemize}

\textbf{3. Que risque-t-on sans convention de nommage ?}\\
Un nommage improvisé rend l'environnement \textbf{illisible et difficile à exploiter} : impossible de savoir
d'un coup d'œil à quoi sert une ressource ni à quel environnement elle appartient. Conséquences : erreurs de
manipulation, audit et analyse de coûts compliqués, automatisation fragile, et \textbf{collisions sur les noms
globaux uniques} (Storage, SQL) faute de suffixe.

\begin{etatbox}
\begin{itemize}
  \item Resource Group \code{rg-shopeasy-dev} créé en \code{swedencentral}, \code{provisioningState: Succeeded}.
  \item 5 tags de gouvernance appliqués ; convention de nommage figée (suffixe \code{ls01}).
  \item \textbf{Prêt pour l'Atelier 5 — création du réseau (VNet + subnets).}
\end{itemize}
\end{etatbox}


% #####################################################################
%  ATELIER 5
% #####################################################################
\section{Atelier 5 — Création du réseau Azure}

\begin{fiche}
\textbf{\eTarget~Objectif} — créer un réseau virtuel segmenté pour séparer les composants web et données.

\textbf{\ePackage~Livrable attendu} — capture du Virtual Network et des subnets, avec commentaire sur la segmentation.

\textbf{\eGlobe~Région} — \code{swedencentral} (Stockholm), seule région autorisée où les VM sont déployables.
\end{fiche}

\subsection{Plan d'adressage}
\begin{ltable}{Q[l,0.9] Q[l,0.6] Q[c,0.5] X[l,1]}
Élément & Plage CIDR & Adresses & Rôle \\
\textbf{VNet} \code{vnet-shopeasy-dev} & \code{10.10.0.0/16} & 65\,536 & Réseau global de ShopEasy \\
\code{snet-web} & \code{10.10.1.0/24} & 251* & Machines virtuelles web \\
\code{snet-data} & \code{10.10.2.0/24} & 251* & Services de données / endpoints privés \\
\code{snet-admin} & \code{10.10.3.0/24} & 251* & Accès d'administration / bastion (optionnel) \\
\end{ltable}
{\small * Azure réserve \textbf{5 adresses par subnet} (la 1\textsuperscript{re}, les 3 suivantes et la dernière) \arr{} $256-5 = 251$ utilisables.}

\subsection{Création via Azure CLI}
\begin{codebox}{bash}
\begin{Verbatim}
# 1) VNet + premier subnet (web)
az network vnet create \
  --resource-group rg-shopeasy-dev --name vnet-shopeasy-dev \
  --location swedencentral --address-prefix 10.10.0.0/16 \
  --subnet-name snet-web --subnet-prefix 10.10.1.0/24

# 2) Subnet data
az network vnet subnet create -g rg-shopeasy-dev --vnet-name vnet-shopeasy-dev \
  --name snet-data --address-prefixes 10.10.2.0/24

# 3) Subnet admin
az network vnet subnet create -g rg-shopeasy-dev --vnet-name vnet-shopeasy-dev \
  --name snet-admin --address-prefixes 10.10.3.0/24
\end{Verbatim}
\end{codebox}

\subsection{Vérification (sortie CLI)}
\begin{minipage}[t]{0.46\linewidth}
\begin{ltable}{Q[l] Q[c]}
VNet & État \\
\code{vnet-shopeasy-dev} & Succeeded \\
\end{ltable}
\end{minipage}\hfill
\begin{minipage}[t]{0.5\linewidth}
\begin{ltable}{Q[l] Q[l] Q[c]}
Subnet & CIDR & État \\
\code{snet-web} & \code{10.10.1.0/24} & Succeeded \\
\code{snet-data} & \code{10.10.2.0/24} & Succeeded \\
\code{snet-admin} & \code{10.10.3.0/24} & Succeeded \\
\end{ltable}
\end{minipage}

\screenshot{vnet-overview.png}{Vue d'ensemble du VNet \code{vnet-shopeasy-dev} (\code{10.10.0.0/16}, Sweden Central).}
\screenshot{vnet-subnets.png}{Les 3 subnets : \code{snet-web}, \code{snet-data}, \code{snet-admin}.}

\subsection{Pourquoi une seule grande plage ne suffit pas ? (segmentation)}
Mettre toutes les ressources dans \textbf{un seul grand subnet « à plat »} fonctionne techniquement, mais
c'est une \textbf{mauvaise pratique} professionnelle :
\begin{itemize}
  \item \textbf{Sécurité (cloisonnement)} : sans segmentation, une VM web compromise atteint \textbf{directement la base}. En séparant \code{snet-web} et \code{snet-data}, on applique des NSG différents et on \textbf{limite les mouvements latéraux}.
  \item \textbf{Moindre privilège réseau} : chaque subnet n'autorise que les flux nécessaires (la base n'accepte 1433 que depuis le subnet web).
  \item \textbf{Lisibilité \& exploitation} : un découpage par fonction (web / données / admin) rend le réseau auditable.
  \item \textbf{Gouvernance \& évolutivité} : route tables, Private Endpoints ou bastion associés à un subnet précis ; plages réservées pour la croissance.
  \item \textbf{Conformité} : isoler les données sensibles facilite le respect du RGPD.
\end{itemize}

\begin{tipbox}[En résumé]
\textbf{Segmenter, c'est appliquer la défense en profondeur au niveau réseau.} La grande plage \code{/16}
sert d'espace global cohérent, et les \code{/24} matérialisent les \textbf{zones de confiance} distinctes.
\end{tipbox}

\begin{etatbox}
\begin{itemize}
  \item VNet \code{vnet-shopeasy-dev} (\code{10.10.0.0/16}) créé en \code{swedencentral}.
  \item 3 subnets \code{Succeeded} : web (10.10.1.0/24), data (10.10.2.0/24), admin (10.10.3.0/24).
  \item \textbf{Prêt pour l'Atelier 6 — filtrage réseau (NSG).}
\end{itemize}
\end{etatbox}


% #####################################################################
%  ATELIER 6
% #####################################################################
\section{Atelier 6 — Filtrage réseau avec Network Security Groups}

\begin{fiche}
\textbf{\eTarget~Objectif} — protéger les flux réseau en limitant les ouvertures inutiles.

\textbf{\ePackage~Livrable attendu} — capture des NSG et des règles, avec justification des ouvertures autorisées.
\end{fiche}

\subsection{Règles de filtrage retenues}
\begin{tblr}{
  width=\linewidth, colspec={Q[l]Q[l]Q[c]Q[c]Q[c]Q[c]Q[l]X[l]},
  row{1}={bg=azuredark,fg=white,font=\sffamily\bfseries\scriptsize},
  row{2-Z}={font=\scriptsize}, row{even}={bg=rowalt},
  hline{1,Z}={1.1pt,azuredark}, hline{2}={0.5pt,azuredark!45},
  colsep=4pt, rowsep=3.2pt, stretch=1.15,
}
NSG & Règle & Prio & Sens & Accès & Port & Source & Justification \\
\code{nsg-web} & \code{allow-http} & 100 & In & Allow & 80 & Internet & Accès HTTP à l'application de test \\
\code{nsg-web} & \code{allow-https} & 110 & In & Allow & 443 & Internet & Accès HTTPS (cible production) \\
\code{nsg-web} & \code{allow-ssh-my-ip} & 120 & In & Allow & 22 & \code{<IP>/32} & Administration Linux limitée à l'IP admin \\
\code{nsg-data} & \code{allow-sql-from-web} & 100 & In & Allow & 1433 & \code{10.10.1.0/24} & Accès SQL uniquement depuis le subnet web \\
\code{nsg-data} & \code{deny-sql-from-vnet} & 200 & In & \textbf{Deny} & 1433 & \code{VNet} & Bloque 1433 depuis le reste du VNet \\
\end{tblr}

\begin{notebox}[Confidentialité]
L'IP publique réelle de l'administrateur est utilisée dans la règle déployée, mais \textbf{masquée ici}
(\code{<IP>/32}) pour ne pas la divulguer.
\end{notebox}

\subsubsection{Pourquoi la règle \code{deny-sql-from-vnet} ?}
Un NSG applique des \textbf{règles par défaut} invisibles, dont \code{AllowVnetInBound} (priorité 65000) qui
\textbf{autorise tout le trafic interne au VNet}. Sans action, \code{snet-admin} pourrait joindre la base sur
le 1433. La règle \code{allow-sql-from-web} (prio 100) autorise le subnet web, et \code{deny-sql-from-vnet}
(prio 200) \textbf{refuse explicitement} le 1433 pour le reste du VNet \arr{} un vrai \textit{« uniquement depuis
le web »} (moindre privilège réseau).

\subsection{Création via Azure CLI}
\begin{codebox}{bash}
\begin{Verbatim}
# ---- NSG WEB ----
az network nsg create -g rg-shopeasy-dev -n nsg-web --location swedencentral
az network nsg rule create -g rg-shopeasy-dev --nsg-name nsg-web \
  --name allow-http --priority 100 --access Allow --direction Inbound \
  --protocol Tcp --source-address-prefixes Internet --destination-port-ranges 80
az network nsg rule create -g rg-shopeasy-dev --nsg-name nsg-web \
  --name allow-https --priority 110 --access Allow --direction Inbound \
  --protocol Tcp --source-address-prefixes Internet --destination-port-ranges 443
az network nsg rule create -g rg-shopeasy-dev --nsg-name nsg-web \
  --name allow-ssh-my-ip --priority 120 --access Allow --direction Inbound \
  --protocol Tcp --source-address-prefixes <VOTRE_IP>/32 --destination-port-ranges 22

# ---- NSG DATA ----
az network nsg create -g rg-shopeasy-dev -n nsg-data --location swedencentral
az network nsg rule create -g rg-shopeasy-dev --nsg-name nsg-data \
  --name allow-sql-from-web --priority 100 --access Allow --direction Inbound \
  --protocol Tcp --source-address-prefixes 10.10.1.0/24 --destination-port-ranges 1433
az network nsg rule create -g rg-shopeasy-dev --nsg-name nsg-data \
  --name deny-sql-from-vnet --priority 200 --access Deny --direction Inbound \
  --protocol Tcp --source-address-prefixes VirtualNetwork --destination-port-ranges 1433

# ---- ASSOCIATION AUX SUBNETS ----
az network vnet subnet update -g rg-shopeasy-dev --vnet-name vnet-shopeasy-dev \
  --name snet-web --network-security-group nsg-web
az network vnet subnet update -g rg-shopeasy-dev --vnet-name vnet-shopeasy-dev \
  --name snet-data --network-security-group nsg-data
\end{Verbatim}
\end{codebox}

\subsection{Vérification — associations subnet \texorpdfstring{$\rightarrow$}{->} NSG}
\begin{ltable}{Q[l] Q[l] Q[l]}
Subnet & CIDR & NSG associé \\
\code{snet-web} & \code{10.10.1.0/24} & \code{nsg-web} \\
\code{snet-data} & \code{10.10.2.0/24} & \code{nsg-data} \\
\code{snet-admin} & \code{10.10.3.0/24} & aucun (optionnel) \\
\end{ltable}

\screenshot{nsg-web-rules.png}{Règles entrantes de \code{nsg-web} (80/443 depuis Internet, 22 depuis l'IP admin) + règles par défaut.}
\screenshot{nsg-data-rules.png}{Règles entrantes de \code{nsg-data} : \code{allow-sql-from-web} (1433) et \code{deny-sql-from-vnet}.}

\subsection{Questions de sécurité}

\textbf{1. Pourquoi limiter SSH à une adresse IP précise ?}\\
Ouvrir le port 22 à \code{0.0.0.0/0} l'expose en permanence aux \textbf{scans automatisés} et aux
\textbf{attaques par force brute}. En restreignant la source à l'IP de l'administrateur, la surface d'attaque
passe de « tout Internet » à « une IP ». En production, on supprime l'exposition publique du SSH au profit
d'\textbf{Azure Bastion} + MFA.

\textbf{2. Pourquoi ne pas exposer le port SQL (1433) sur Internet ?}\\
La base contient les \textbf{données sensibles}. L'exposer en ferait une \textbf{cible privilégiée} : force
brute, exploitation de vulnérabilités, exfiltration. Le 1433 ne doit être joignable que depuis la couche
applicative (subnet web), idéalement via \textbf{Private Endpoint}.

\textbf{3. Différence entre une règle entrante et une règle sortante ?}
\begin{itemize}
  \item \textbf{Entrante (Inbound)} : filtre le trafic qui \textbf{arrive} vers la ressource (requêtes vers le web, ou du web vers la base).
  \item \textbf{Sortante (Outbound)} : filtre le trafic qui \textbf{part} de la ressource (appel à un service externe).
\end{itemize}
Un NSG évalue les deux sens \textbf{séparément}, par priorité. Par défaut Azure autorise tout l'Outbound et
restreint l'Inbound. Restreindre l'Outbound limite l'\textbf{exfiltration}.

\textbf{4. Comment cette configuration devrait évoluer en production ?}
\begin{itemize}
  \item \textbf{HTTPS obligatoire} + redirection 80\arr{}443, certificats ; \textbf{Application Gateway + WAF}.
  \item \textbf{Suppression du SSH public} \arr{} Azure Bastion, VM sans IP publique, accès just-in-time.
  \item \textbf{Private Endpoint} pour Azure SQL \arr{} plus aucune ouverture 1433.
  \item \textbf{NSG Flow Logs} + Azure Monitor ; règles Outbound restrictives ; revue périodique + MFA.
\end{itemize}

\begin{etatbox}
\begin{itemize}
  \item \code{nsg-web} (3 règles) associé à \code{snet-web} ; \code{nsg-data} (2 règles) à \code{snet-data}.
  \item Base non joignable hors du subnet web ; SSH limité à l'IP admin.
  \item \textbf{Prêt pour l'Atelier 7 — déploiement des 2 VM web.}
\end{itemize}
\end{etatbox}


% #####################################################################
%  ATELIER 7
% #####################################################################
\section{Atelier 7 — Déploiement de machines virtuelles web}

\begin{fiche}
\textbf{\eTarget~Objectif} — déployer deux serveurs web simples dans Azure pour représenter la couche applicative.

\textbf{\ePackage~Livrable attendu} — captures des deux VM, de la page web de test et des commandes de vérification.
\end{fiche}

\begin{notebox}[Adaptation imposée par Azure for Students (région + taille)]
Le TP demande \code{Standard\_B1s} en \code{francecentral}. \textbf{Aucun des deux n'est utilisable} sur cet
abonnement. Après analyse :
\begin{itemize}
  \item \code{francecentral} bloquée par policy ; \code{germanywestcentral} : tailles x86 \textit{NotAvailableForSubscription}, ARM à quota 0.
  \item \code{Standard\_B1s} \textit{NotAvailable} dans les 5 régions autorisées ; providers \code{Compute/Storage/Sql} étaient \textbf{NotRegistered} (enregistrés \arr{} quota débloqué).
\end{itemize}
\textbf{Solution retenue} (appuyée par la doc Microsoft) : région \textbf{\code{swedencentral}} + taille
\textbf{\code{Standard\_B2ats\_v2}} (2 vCPU, 1 Go), \textbf{taille officiellement recommandée pour Azure for
Students} (famille \code{Basv2}, quota 10 vCPU).
\end{notebox}

\subsection{Création des VM (Azure CLI)}
\begin{codebox}{bash}
\begin{Verbatim}
az vm create \
  --resource-group rg-shopeasy-dev --name vm-web-01 \
  --image Ubuntu2204 --size Standard_B2ats_v2 --location swedencentral \
  --vnet-name vnet-shopeasy-dev --subnet snet-web --nsg nsg-web \
  --admin-username azureuser --generate-ssh-keys \
  --public-ip-sku Standard --storage-sku Standard_LRS

# vm-web-02 : idem avec --name vm-web-02
\end{Verbatim}
\end{codebox}

\begin{ltable}{Q[l] Q[l] Q[l] Q[l]}
VM & IP publique & IP privée & Taille \\
\code{vm-web-01} & \code{20.240.255.137} & \code{10.10.1.4} & \code{Standard\_B2ats\_v2} \\
\code{vm-web-02} & \code{4.223.107.224} & \code{10.10.1.5} & \code{Standard\_B2ats\_v2} \\
\end{ltable}

\subsection{Installation de Nginx (SSH)}
Commandes lancées sur chaque VM (connexion SSH avec la clé générée \code{\textasciitilde/.ssh/id\_rsa}) :
\begin{codebox}{bash}
\begin{Verbatim}
ssh azureuser@<IP_PUBLIQUE>
sudo apt update
sudo apt install -y nginx
echo "ShopEasy - VM Web 0X" | sudo tee /var/www/html/index.html
sudo systemctl enable nginx
sudo systemctl restart nginx
systemctl status nginx --no-pager
\end{Verbatim}
\end{codebox}

Sortie réelle de \code{systemctl status} sur \code{vm-web-01} (service \textbf{actif}) :
\begin{codebox}{console}
\begin{Verbatim}
########## sudo apt install -y nginx ##########
0 upgraded, 20 newly installed, 0 to remove and 7 not upgraded.
...
Setting up nginx-common (1.18.0-6ubuntu14.16) ...
Created symlink /etc/systemd/system/multi-user.target.wants/nginx.service.
Setting up nginx (1.18.0-6ubuntu14.16) ...

########## page index.html ##########
ShopEasy - VM Web 01

########## systemctl status nginx ##########
* nginx.service - A high performance web server and a reverse proxy server
     Loaded: loaded (/lib/systemd/system/nginx.service; enabled; vendor preset: enabled)
     Active: active (running) since Wed 2026-06-24 10:55:29 UTC; 6ms ago
   Main PID: 4297 (nginx)
      Tasks: 3 (limit: 1064)
     Memory: 3.3M
     CGroup: /system.slice/nginx.service
             |-4297 "nginx: master process /usr/sbin/nginx ..."
             |-4298 "nginx: worker process"
             `-4299 "nginx: worker process"
\end{Verbatim}
\end{codebox}
\textit{La page de \code{vm-web-02} affiche « ShopEasy - VM Web 02 » ; son service Nginx est également}
\textbf{\textit{active (running)}}\textit{ et }\textbf{\textit{enabled}}.

\subsection{Vérifications (commandes CLI)}
Les deux VM dans le Resource Group, état \code{VM running} :
\begin{ltable}{Q[l] Q[l] Q[l] Q[l] Q[c]}
Nom & Taille & IP publique & IP privée & État \\
\code{vm-web-01} & \code{Standard\_B2ats\_v2} & \code{20.240.255.137} & \code{10.10.1.4} & VM running \\
\code{vm-web-02} & \code{Standard\_B2ats\_v2} & \code{4.223.107.224} & \code{10.10.1.5} & VM running \\
\end{ltable}

Test HTTP de chaque page depuis le poste admin :
\begin{codebox}{console}
\begin{Verbatim}
$ curl http://20.240.255.137
ShopEasy - VM Web 01
$ curl http://4.223.107.224
ShopEasy - VM Web 02
\end{Verbatim}
\end{codebox}

\textbf{Réponses aux questions de vérification.} \textbf{(1)} Les deux VM apparaissent dans le Resource Group
(état \code{VM running}). \textbf{(2)} Nginx est \code{active (running)} + \code{enabled} sur les deux.
\textbf{(3)} Les pages sont accessibles (\code{curl}, NSG \code{allow-http} 80). \textbf{(4)} Les deux pages
identifient la VM appelée (« VM Web 01 » / « VM Web 02 »), ce qui servira à vérifier la répartition de charge
à l'Atelier 8.

\screenshot{vm-web-01-page.png}{Page servie par \code{vm-web-01} (\code{http://20.240.255.137}).}
\screenshot{vm-web-02-page.png}{Page servie par \code{vm-web-02} (\code{http://4.223.107.224}).}
\screenshot{vms-in-rg.png}{Les deux VM dans le Resource Group \code{rg-shopeasy-dev}.}

\begin{etatbox}
\begin{itemize}
  \item 2 VM \code{Standard\_B2ats\_v2} (Ubuntu 22.04) dans \code{snet-web}, Nginx actif, pages distinctes.
  \item \textbf{Prêt pour l'Atelier 8 — répartiteur de charge devant les 2 VM.}
\end{itemize}
\end{etatbox}


% #####################################################################
%  ATELIER 8
% #####################################################################
\section{Atelier 8 — Répartition de charge}

\begin{fiche}
\textbf{\eTarget~Objectif} — comprendre le rôle d'un répartiteur de charge et l'intégrer dans l'architecture cible.

\textbf{\ePackage~Livrable attendu} — capture du répartiteur, du backend pool, de la sonde de santé et du test navigateur.

\textbf{\eGlobe~Région} — \texttt{swedencentral} \quad\textbullet\quad \textbf{Load Balancer} \texttt{lb-shopeasy} (SKU Standard, public).
\end{fiche}

\subsection{Choix technique}
\begin{ltable}{Q[l,1.0] X[l,1.5] X[l,1.5]}
Critère & Azure Load Balancer \textbf{(retenu)} & Application Gateway \\
Niveau & Couche 4 (TCP/UDP) & Couche 7 (HTTP/HTTPS) \\
Usage & Répartition réseau simple et efficace & Routage applicatif, terminaison TLS, \textbf{WAF} \\
Complexité & Plus simple & Plus riche, plus complexe \\
Cas TP & \eCheck~\textbf{Suffisant} pour une première architecture web & Cible production (sécurité applicative) \\
\end{ltable}

\noindent On retient \textbf{Azure Load Balancer (Standard, public)}, conforme à l'esprit du TP.
\textit{Application Gateway + WAF est la cible production} (cf. variante de durcissement à l'Atelier~3).

\subsection{Création via Azure CLI}
\begin{codebox}{bash}
\begin{Verbatim}
# 1) IP publique frontale (Standard, statique)
az network public-ip create -g rg-shopeasy-dev -n pip-lb-shopeasy \
  --sku Standard --location swedencentral --allocation-method Static

# 2) Load Balancer + frontend + backend pool
az network lb create -g rg-shopeasy-dev -n lb-shopeasy \
  --sku Standard --location swedencentral \
  --public-ip-address pip-lb-shopeasy \
  --frontend-ip-name feip-shopeasy \
  --backend-pool-name bepool-web

# 3) Sonde de sante HTTP sur le port 80 (chemin /)
az network lb probe create -g rg-shopeasy-dev --lb-name lb-shopeasy \
  --name probe-http --protocol Http --port 80 --path / --interval 5

# 4) Regle de repartition HTTP (80 -> 80) liee a la sonde
az network lb rule create -g rg-shopeasy-dev --lb-name lb-shopeasy \
  --name rule-http --protocol Tcp --frontend-port 80 --backend-port 80 \
  --frontend-ip-name feip-shopeasy --backend-pool-name bepool-web \
  --probe-name probe-http

# 5) Rattacher les 2 VM au backend pool (via leurs ipconfig)
az network nic ip-config address-pool add -g rg-shopeasy-dev \
  --nic-name vm-web-01VMNic --ip-config-name ipconfigvm-web-01 \
  --lb-name lb-shopeasy --address-pool bepool-web
az network nic ip-config address-pool add -g rg-shopeasy-dev \
  --nic-name vm-web-02VMNic --ip-config-name ipconfigvm-web-02 \
  --lb-name lb-shopeasy --address-pool bepool-web
\end{Verbatim}
\end{codebox}

\subsection{Vérification (sorties CLI réelles)}
\noindent IP publique du Load Balancer : \texttt{20.91.244.18}. Membres du backend pool : les deux interfaces
\texttt{vm-web-01VMNic} et \texttt{vm-web-02VMNic}.

\smallskip
\begin{minipage}[t]{0.48\linewidth}
\begin{ltable}{Q[l] Q[c] Q[c] Q[c]}
Sonde & Proto & Port & Chemin \\
\texttt{probe-http} & Http & 80 & / \\
\end{ltable}
\end{minipage}\hfill
\begin{minipage}[t]{0.48\linewidth}
\begin{ltable}{Q[l] Q[c] Q[c] Q[c]}
Règle & Proto & Front & Back \\
\texttt{rule-http} & Tcp & 80 & 80 \\
\end{ltable}
\end{minipage}

\subsubsection{Test de répartition (preuve de bascule)} 14 requêtes successives sur l'IP publique du
Load Balancer, \textbf{réparties 50/50} entre les deux VM :
\begin{codebox}{console}
\begin{Verbatim}
$ for i in $(seq 1 14); do curl -s http://20.91.244.18; done | sort | uniq -c
   7 ShopEasy - VM Web 01
   7 ShopEasy - VM Web 02
\end{Verbatim}
\end{codebox}
\begin{tipbox}[Suppression du SPOF]
Le Load Balancer distribue bien le trafic sur les \textbf{deux} serveurs web : c'est la suppression du
\textbf{point de défaillance unique} applicatif identifié à l'Atelier~1.
\end{tipbox}

\subsection{Captures}
\screenshot{lb-overview.png}{Vue d'ensemble du Load Balancer — IP publique \texttt{20.91.244.18}, SKU Standard.}
\screenshot{lb-backend-pool.png}{Backend pool \texttt{bepool-web} : \texttt{vm-web-01} (10.10.1.4) et \texttt{vm-web-02} (10.10.1.5).}
\screenshot{lb-health-probe.png}{Sonde de santé \texttt{probe-http} (HTTP 80, chemin /), utilisée par \texttt{rule-http}.}
\screenshot{lb-test-browser.png}{Test navigateur : la même IP sert « VM Web 01 » et « VM Web 02 » (bascule).}

\subsection{Analyse (questions du TP)}
\textbf{1. Que se passe-t-il si une VM devient indisponible ?}\\
La \textbf{sonde de santé} cesse de recevoir une réponse HTTP 200 sur \texttt{/} pour cette VM et la
\textbf{retire automatiquement du backend pool}. Le Load Balancer route alors 100\,\% du trafic vers la VM
saine restante : le service \textbf{reste disponible} (capacité réduite, mais pas d'interruption). Dès que
la VM répond à nouveau à la sonde, elle est \textbf{réintégrée} automatiquement.

\smallskip\textbf{2. Pourquoi une sonde de santé est-elle nécessaire ?}\\
Sans sonde, le LB continuerait d'envoyer du trafic vers une instance morte $\rightarrow$ \textbf{erreurs
pour une partie des utilisateurs}. La sonde vérifie en continu que chaque backend répond réellement, et
\textbf{ne route que vers les instances saines}. C'est elle qui rend la redondance \textit{effective}.

\smallskip\textbf{3. Quelles limites subsistent malgré le répartiteur de charge ?}
\begin{itemize}
  \item \textbf{Pas de protection inter-zone/région} : les 2 VM sont dans la même région, \textbf{sans Availability Zone} (déploiement zonal restreint sur Azure for Students) $\rightarrow$ une panne de datacenter peut toucher les deux.
  \item \textbf{Pas d'autoscaling} : le nombre de VM est fixe, l'ajout est manuel.
  \item \textbf{Couche 4 seulement} : pas de terminaison \textbf{TLS} ni de \textbf{WAF} ($\rightarrow$ Application Gateway en production).
  \item \textbf{Couche données non redondée} ici (une seule base).
  \item L'application doit rester \textbf{sans état} pour que la répartition soit transparente.
\end{itemize}

\smallskip\textbf{4. Pourquoi déployer les VM dans plusieurs zones améliore-t-il la disponibilité ?}\\
Les \textbf{Availability Zones} sont des datacenters \textbf{physiquement séparés} (alimentation,
refroidissement, réseau indépendants) au sein d'une même région. En plaçant \texttt{vm-web-01} et
\texttt{vm-web-02} dans \textbf{deux zones distinctes}, une panne touchant un datacenter n'affecte
\textbf{qu'une seule} VM. \textit{(Point repris à l'Atelier~13.)}

\begin{etatbox}
\begin{itemize}
  \item Load Balancer public \texttt{lb-shopeasy} (\texttt{20.91.244.18}) devant les 2 VM, sonde HTTP/80, répartition 50/50 vérifiée.
  \item Point de défaillance unique applicatif \textbf{levé}.
  \item \textbf{Prêt pour l'Atelier~9 — stockage documentaire (Storage Account).}
\end{itemize}
\end{etatbox}


% #####################################################################
%  ATELIER 9
% #####################################################################
\section{Atelier 9 — Stockage documentaire avec Azure Storage}

\begin{fiche}
\textbf{\eTarget~Objectif} — remplacer le stockage local des documents par un service de stockage objet.

\textbf{\ePackage~Livrable attendu} — capture du Storage Account, des conteneurs, du téléversement de fichiers et réponses aux questions.

\textbf{\eGlobe~Région} — \code{swedencentral} \quad\textbullet\quad \textbf{Storage Account} \code{stshopeasyls01} (suffixe \code{ls01}).
\end{fiche}

\subsection{Création du Storage Account (Azure CLI)}
\begin{codebox}{bash}
\begin{Verbatim}
az storage account create \
  --resource-group rg-shopeasy-dev --name stshopeasyls01 \
  --location swedencentral --sku Standard_LRS --kind StorageV2 \
  --https-only true --min-tls-version TLS1_2 \
  --allow-blob-public-access false \
  --tags projet=shopeasy environnement=dev module=cloud
\end{Verbatim}
\end{codebox}
{\small Résultat : \code{provisioningState: Succeeded}, \code{Https: true}, \code{TLS: TLS1\_2}, \code{AccesPublicBlob: false}.}

\subsection{Sécurité et résilience des données}
\begin{ltable}{Q[l,0.7] Q[l,0.6] X[l,1]}
Propriété & Valeur & Apport \\
\textbf{Chiffrement au repos} & activé (Microsoft) & Données chiffrées par défaut (clés gérées Microsoft) \\
\textbf{HTTPS obligatoire} & \code{true} & Données chiffrées en transit \\
\textbf{TLS minimum} & \code{TLS1\_2} & Pas de protocole obsolète \\
\textbf{Accès public blob} & \code{false} & Impossible de rendre un conteneur public par erreur \\
\textbf{Versioning des blobs} & \code{true} & Conserve chaque version (anti-écrasement / anti-suppression) \\
\textbf{Soft-delete blob} & 7 jours & Récupération d'un blob supprimé \\
\textbf{Soft-delete conteneur} & 7 jours & Récupération d'un conteneur supprimé \\
\end{ltable}
{\small Le chiffrement étant \textbf{activé par défaut} sur StorageV2, le point « activer le chiffrement » du TP est satisfait nativement.}

\subsection{Conteneurs et téléversement}
\begin{codebox}{bash}
\begin{Verbatim}
# Conteneurs privés
for C in factures clients archives; do
  az storage container create --account-name stshopeasyls01 --name "$C" --public-access off
done

# Authentification data-plane via variables d'environnement (cle non exposee)
export AZURE_STORAGE_ACCOUNT=stshopeasyls01
export AZURE_STORAGE_KEY=<cle du compte>

az storage blob upload -c factures -f facture-2026-001.txt -n facture-2026-001.txt
az storage blob upload -c clients  -f client-dupont.txt    -n client-dupont.txt
# ... (5 fichiers de test au total)
\end{Verbatim}
\end{codebox}

Contenu vérifié (\code{az storage blob list}) :
\begin{ltable}{Q[l,0.5] Q[l,1] Q[c,0.4]}
Conteneur & Blob & Taille (o) \\
\code{factures} & \code{facture-2026-001.txt} & 69 \\
\code{factures} & \code{facture-2026-002.txt} & 69 \\
\code{clients} & \code{client-dupont.txt} & 72 \\
\code{clients} & \code{client-martin.txt} & 70 \\
\code{archives} & \code{archive-commandes-2025.txt} & 62 \\
\end{ltable}

\subsection{Politique de cycle de vie (FinOps) — bonus appliqué}
\begin{codebox}{json}
\begin{Verbatim}
{ "rules": [{ "enabled": true, "name": "archives-tiering", "type": "Lifecycle",
  "definition": {
    "filters": { "blobTypes": ["blockBlob"], "prefixMatch": ["archives/"] },
    "actions": { "baseBlob": {
      "tierToCool":    { "daysAfterModificationGreaterThan": 30 },
      "tierToArchive": { "daysAfterModificationGreaterThan": 90 },
      "delete":        { "daysAfterModificationGreaterThan": 365 } } } } }] }
\end{Verbatim}
\end{codebox}
\begin{tipbox}[Cycle de vie des archives]
Règle \code{archives-tiering} active sur le préfixe \code{archives/} : \textbf{Hot \arr{} Cool (30 j) \arr{} Archive (90 j) \arr{} suppression (365 j)}. Réduit les coûts sans intervention manuelle.
\end{tipbox}

\screenshot{storage-account.png}{Storage Account \code{stshopeasyls01} : HTTPS-only, TLS 1.2, accès public désactivé, versioning et soft-delete activés.}
\screenshot{storage-containers.png}{Les 3 conteneurs privés (\code{factures}, \code{clients}, \code{archives}).}
\screenshot{storage-blobs.png}{Conteneur \code{factures} avec ses 2 fichiers téléversés.}

\subsection{Questions du TP}

\textbf{Comment ce service remplace-t-il le stockage local du serveur historique ?}\\
À l'origine, les documents étaient dans un \textbf{répertoire local} de la VM unique : couplés au serveur,
perdus si le disque tombe, non partagés et non sauvegardés. Le \textbf{Blob Storage} externalise ces fichiers
dans un service \textbf{durable, redondé, chiffré et indépendant des VM}, accessible via API/SDK. On supprime
ainsi le couplage et le risque de perte lié au serveur.

\textbf{1. Pourquoi un stockage objet plutôt qu'un disque local de VM ?}
\begin{itemize}
  \item \textbf{Durabilité \& redondance} gérées par Azure, vs un disque lié à une seule VM.
  \item \textbf{Découplage} : les fichiers survivent à la destruction/au redéploiement des VM.
  \item \textbf{Scalabilité} quasi illimitée, \textbf{paiement à l'usage}, \textbf{accès partagé} par plusieurs instances.
  \item \textbf{Fonctions intégrées} : versioning, soft-delete, cycle de vie, tags, journaux.
\end{itemize}

\textbf{2. Pourquoi éviter les conteneurs publics par défaut ?}\\
Un conteneur public expose les documents à \textbf{tout Internet} sans authentification \arr{} \textbf{fuite de
données} clients (RGPD). Ici, l'accès public est \textbf{bloqué au niveau du compte}, ce qui empêche même une
erreur de configuration.

\textbf{3. Quel intérêt a le versioning ?}\\
Il conserve \textbf{chaque version} d'un blob ; on peut \textbf{restaurer} un document écrasé ou corrompu.
Couplé au soft-delete, il protège contre les suppressions accidentelles ou malveillantes (ransomware).

\textbf{4. Quelle politique de cycle de vie pour les archives ?}\\
Celle appliquée ci-dessus : \textbf{Cool à 30 j}, \textbf{Archive à 90 j}, \textbf{suppression à 365 j} —
réduction des coûts automatique tout en respectant une durée de rétention.

\begin{etatbox}
\begin{itemize}
  \item Storage Account \code{stshopeasyls01} : HTTPS-only, TLS 1.2, accès public bloqué, chiffré, versioning + soft-delete.
  \item 3 conteneurs privés avec 5 fichiers de test ; lifecycle \code{archives-tiering} active.
  \item \textbf{Prêt pour l'Atelier 10 — base de données managée (Azure SQL Database).}
\end{itemize}
\end{etatbox}


% #####################################################################
%  ATELIER 10
% #####################################################################
\section{Atelier 10 — Base de données managée (Azure SQL Database)}

\begin{fiche}
\textbf{\eTarget~Objectif} — comparer une base installée sur serveur avec une base managée Azure SQL Database.

\textbf{\ePackage~Livrable attendu} — capture de la base créée et tableau comparatif complété.

\textbf{\eGlobe~Région} — \code{swedencentral} \quad\textbullet\quad serveur \code{sql-shopeasy-ls01} \quad\textbullet\quad base \code{sqldb-shopeasy}.
\end{fiche}

\subsection{Architecture attendue}
La base de données \textbf{ne doit pas être exposée librement sur Internet}. En production, on privilégie
\textbf{règles réseau limitées}, \textbf{Private Endpoint}, \textbf{identités managées} (Entra ID) et
politiques de sécurité adaptées. Pour ce TP, le serveur est créé avec un \textbf{pare-feu fermé} (aucune règle
= aucun accès entrant).

\subsection{Création (Azure CLI)}
\begin{codebox}{bash}
\begin{Verbatim}
# Serveur logique SQL (mot de passe admin fort genere)
az sql server create -g rg-shopeasy-dev --name sql-shopeasy-ls01 \
  --location swedencentral --admin-user shopeasyadmin \
  --admin-password '<MOT_DE_PASSE_FORT>'

# Base managee, niveau Basic (test)
az sql db create -g rg-shopeasy-dev --server sql-shopeasy-ls01 \
  --name sqldb-shopeasy --service-objective Basic --backup-storage-redundancy Local
\end{Verbatim}
\end{codebox}
\begin{notebox}[Identifiants]
À conserver hors dépôt Git : admin \code{shopeasyadmin}, mot de passe \textbf{fort généré} (non publié).
Serveur : \code{sql-shopeasy-ls01.database.windows.net}.
\end{notebox}

{\small Sorties CLI : serveur \code{Ready} ; base \code{sqldb-shopeasy} — \code{Tier: Basic}, \code{DTU: 5}, \code{MaxSize: 2 Go}, \code{Etat: Online}.}

\subsection{Sécurité — base non exposée}
\begin{codebox}{console}
\begin{Verbatim}
=== Regles de pare-feu du serveur SQL ===
  Nombre de regles : 0   (0 regle = aucune IP autorisee = base NON exposee a Internet)
=== Acces reseau public + TLS ===
  { "AccesPublic": "Enabled", "TLS_min": "1.2" }
\end{Verbatim}
\end{codebox}
\begin{itemize}
  \item \textbf{Aucune règle de pare-feu} \arr{} le point de terminaison public existe mais \textbf{rejette toutes les IP}.
  \item Pour se connecter en test, on ajouterait \textbf{temporairement} une règle limitée à l'IP admin, documentée et supprimée après validation.
  \item \textbf{Cible production} : \code{publicNetworkAccess=Disabled} + \textbf{Private Endpoint} dans \code{snet-data}.
\end{itemize}

\screenshot{sql-database.png}{Base \code{sqldb-shopeasy} : serveur \code{sql-shopeasy-ls01}, tier Basic, statut Online.}

\subsection{Tableau comparatif — SQL sur VM vs Azure SQL Database}
\begin{ltable}{Q[l,0.6] X[l,1] X[l,1]}
Critère & Base sur VM (IaaS) & Azure SQL Database (PaaS) \\
\textbf{Administration OS} & À la charge du client (OS + moteur SQL) & \textbf{Aucune} — Azure gère OS et moteur \\
\textbf{Sauvegardes} & À concevoir et opérer & \textbf{Automatiques} (point-in-time restore) \\
\textbf{Mises à jour} & À planifier et appliquer & \textbf{Automatiques} et transparentes \\
\textbf{Haute disponibilité} & À construire (cluster, réplication) & \textbf{Intégrée} (SLA jusqu'à 99,99 \%) \\
\textbf{Sécurité} & Durcissement à la charge du client & \textbf{TDE} par défaut, pare-feu, Entra ID, audit \\
\textbf{Coût} & VM + licence + disque + exploitation & Paiement du \textbf{service} (DTU/vCore) \\
\textbf{Flexibilité} & Contrôle total (version, OS) & Périmètre base, mais \textbf{scaling rapide} \\
\end{ltable}

\textbf{Synthèse.} Pour ShopEasy, \textbf{Azure SQL Database} supprime l'essentiel de l'administration (OS,
sauvegardes, HA, patchs) identifiée comme un risque à l'Atelier 1, au prix d'un contrôle moindre sur la
plateforme — un compromis pertinent pour une PME sans équipe DBA dédiée.

\begin{etatbox}
\begin{itemize}
  \item Serveur \code{sql-shopeasy-ls01} + base \code{sqldb-shopeasy} (Basic) \textbf{Online}, \textbf{pare-feu fermé}.
  \item Tableau comparatif VM vs PaaS complété.
  \item \textbf{Prêt pour l'Atelier 11 — supervision (Azure Monitor).}
\end{itemize}
\end{etatbox}


% #####################################################################
%  ATELIER 11
% #####################################################################
\section{Atelier 11 — Supervision avec Azure Monitor}

\begin{fiche}
\textbf{\eTarget~Objectif} — mettre en place une première lecture opérationnelle de l'infrastructure.

\textbf{\ePackage~Livrable attendu} — capture des métriques, de l'alerte et tableau d'indicateurs.
\end{fiche}

\subsection{Métriques d'une VM (CPU)}
\begin{codebox}{bash}
\begin{Verbatim}
VMID=$(az vm show -g rg-shopeasy-dev -n vm-web-01 --query id -o tsv)
az monitor metrics list --resource "$VMID" \
  --metric "Percentage CPU" --interval PT5M --aggregation Average Maximum
\end{Verbatim}
\end{codebox}
Sortie réelle (\code{vm-web-01}, \% CPU — VM au repos servant une page statique) :
\begin{ltable}{Q[l] Q[c] Q[c]}
Heure (UTC) & CPU moyen & CPU max \\
11:47 & 0,244 & 0,92 \\
11:52 & 0,168 & 0,22 \\
11:57 & 0,172 & 0,23 \\
12:02 & 0,168 & 0,22 \\
12:07 & 0,168 & 0,23 \\
12:12 & 0,169 & 0,22 \\
\end{ltable}

\subsection{Création d'une alerte CPU}
\begin{codebox}{bash}
\begin{Verbatim}
az monitor metrics alert create --name alert-cpu-vm-web-01 \
  --resource-group rg-shopeasy-dev --scopes "$VMID" \
  --condition "avg Percentage CPU > 80" \
  --window-size 5m --evaluation-frequency 1m --severity 2 \
  --description "CPU moyen de vm-web-01 superieur a 80% pendant 5 minutes"
\end{Verbatim}
\end{codebox}
{\small L'alerte \code{alert-cpu-vm-web-01} (severity 2, fenêtre 5 min, évaluation 1 min) est \textbf{active}. Le provider \code{microsoft.insights} s'est auto-enregistré à cette occasion.}

\subsection{Logs d'activité du Resource Group}
Synthèse des opérations réalisées pendant le TP (50 événements administratifs) :
\begin{ltable}{X[l,1] Q[c,0.3] Q[c,0.2]}
Opération (action) & Statut & Nb \\
\code{Microsoft.Storage/.../listKeys} & Succeeded & 5 \\
\code{Microsoft.Sql/servers/databases/write} & Succeeded & 3 \\
\code{Microsoft.Network/loadBalancers/write} & Succeeded & 3 \\
\code{Microsoft.Sql/servers/write} & Succeeded & 1 \\
\code{Microsoft.Network/networkSecurityGroups/write} & Succeeded & 2 \\
\code{Microsoft.Network/networkInterfaces/write} & Succeeded & 2 \\
\code{Microsoft.Storage/.../managementPolicies/write} & Succeeded & 1 \\
\end{ltable}
{\small Le journal d'activité \textbf{retrace l'historique} de tout ce qui a été fait, par qui et avec quel résultat (traçabilité / audit).}

\screenshot{monitor-vm-metrics.png}{Métriques de \code{vm-web-01} : graphique \code{Percentage CPU} (avec le pic d'installation).}
\screenshot{monitor-alert.png}{Règle d'alerte \code{alert-cpu-vm-web-01} : condition \code{Percentage CPU > 80}, severity 2.}

\subsection{Tableau d'indicateurs}
\begin{ltable}{Q[l,0.55] Q[l,0.55] X[l,1.1]}
Indicateur & Seuil proposé & Pourquoi le surveiller ? \\
\textbf{CPU VM} & $> 80\,\%$ pendant 5 min & Détecter une \textbf{saturation} \arr{} dégradation des temps de réponse \\
\textbf{Disponibilité HTTP} & échec sonde LB & Détecter une \textbf{VM/backend HS} \arr{} indisponibilité du service \\
\textbf{Espace disque} & $> 85\,\%$ utilisé & Éviter le \textbf{disque plein} (OS/app qui plante) \\
\textbf{Échecs de connexion} & $>$ seuil / minute & Détecter erreurs ou \textbf{tentatives d'intrusion} (brute force) \\
\textbf{Coût journalier} & $>$ budget / 30 & Alerter sur un \textbf{dépassement budgétaire} (FinOps) \\
\end{ltable}

\textbf{Trois indicateurs supplémentaires utiles en production :}
\begin{enumerate}
  \item \textbf{Latence applicative (P95)} \arr{} reflète l'expérience utilisateur réelle.
  \item \textbf{Santé des backends du Load Balancer} (Dip availability) \arr{} détecter une bascule.
  \item \textbf{DTU / connexions / deadlocks Azure SQL} \arr{} santé de la couche données.
\end{enumerate}

\begin{etatbox}
\begin{itemize}
  \item Métriques CPU consultées, alerte \code{alert-cpu-vm-web-01} active, logs d'activité analysés.
  \item Tableau d'indicateurs + 3 indicateurs production proposés.
  \item \textbf{Prêt pour l'Atelier 12 — estimation et optimisation des coûts (FinOps).}
\end{itemize}
\end{etatbox}


% #####################################################################
%  ATELIER 12
% #####################################################################
\section{Atelier 12 — Estimation et optimisation des coûts (FinOps)}

\begin{fiche}
\textbf{\eTarget~Objectif} — produire une première estimation budgétaire et identifier des optimisations FinOps simples.

\textbf{\ePackage~Livrable attendu} — une estimation de coût et trois recommandations FinOps.
\end{fiche}

\subsection{Méthode et source des prix}
Le \textbf{coût réel} (Cost Management) n'apparaît qu'après \textbf{24-48 h} sur un abonnement récent. Les prix
ci-dessous sont \textbf{vérifiés en direct} via l'\textbf{API officielle Azure Retail Prices}
(\code{prices.azure.com}, juin 2026), en pay-as-you-go USD, hypothèse \textbf{24/7} (730 h/mois).

\begin{ltable}{X[l,1.2] Q[l,0.9]}
Ressource déployée (vérifiée en CLI) & Prix unitaire officiel \\
2 × VM \code{Standard\_B2ats\_v2} (Linux) & \textbf{0,00972 \$/h} \\
2 × disque OS Standard HDD S4 (32 Go) LRS & \textbf{1,536 \$/mois} \\
1 × Load Balancer Standard (5 règles incluses) & \textbf{0,025 \$/h} (+ 0,005 \$/Go) \\
3 × IP publique Standard static & \textbf{0,005 \$/h} \\
1 × Azure SQL Database Basic (5 DTU) & \textbf{0,161 \$/jour} \\
1 × règle d'alerte Azure Monitor & $\approx$ 0,10 \$/mois \\
Storage Account StorageV2 LRS (faible volume) & négligeable \\
\end{ltable}

\subsection{Tableau de coûts (prix réels, 24/7)}
\begin{ltable}{Q[l,0.7] Q[l,0.7] Q[c,0.4] X[l,1.2]}
Ressource & Calcul & \$/mois & Optimisation possible \\
\textbf{VM web 1} & 0,00972 × 730 & \textbf{7,10} & Désallouer hors usage (\code{az vm deallocate}) \\
\textbf{VM web 2} & 0,00972 × 730 & \textbf{7,10} & Idem ; autoscale en production \\
\textbf{Disques} & 2 × 1,536 & \textbf{3,07} & Taille adaptée ; suppression au nettoyage \\
\textbf{Load Balancer} & 0,025 × 730 & \textbf{18,25} & App Gateway seulement si L7/WAF ; arrêt hors usage \\
\textbf{IP publiques (3)} & 3 × 0,005 × 730 & \textbf{10,95} & Retirer les IP des VM (Bastion) \arr{} 1 seule \\
\textbf{Azure SQL Database} & 0,161 × 30,4 & \textbf{4,90} & Serverless auto-pause / niveau adapté \\
\textbf{Storage Account} & faible volume & \textbf{<\,0,10} & Lifecycle (déjà appliquée), LRS (pas GRS) \\
\textbf{Monitoring} & 1 alerte & \textbf{0,10} & Métriques de plateforme gratuites \\
\textbf{Total} & & \textbf{$\approx$\,51,5} & ($\approx 0$ si tout désalloué/supprimé hors TP) \\
\end{ltable}

\begin{notebox}[Postes les plus coûteux]
Le \textbf{Load Balancer (18,25 \$)}, les \textbf{IP publiques (10,95 \$)} et le \textbf{compute des 2 VM
(14,20 \$)} — le LB et les IP facturent \textbf{même à trafic nul}. Désallouer les VM annule leur coût compute
($\approx$ 14 \$/mois économisés). \textit{Prix vérifiés via l'API Azure Retail Prices (juin 2026).}
\end{notebox}

\subsection{Trois recommandations FinOps}
\begin{enumerate}
  \item \textbf{Arrêter le compute hors usage.} \code{az vm deallocate} stoppe la facturation compute ; en fin de TP, \code{az group delete} supprime d'un coup les postes qui coûtent à vide (LB, IP, base).
  \item \textbf{Réduire les ressources facturées en permanence.} Supprimer les IP publiques des VM (3 \arr{} 1), garder LRS (pas de géo-redondance en dev), conserver la lifecycle policy.
  \item \textbf{Dimensionner et piloter.} Niveaux adaptés (SQL Basic ou Serverless auto-pause), gouvernance via \textbf{tags + budgets/alertes} Cost Management.
\end{enumerate}

\subsection{Questions FinOps du TP}

\textbf{1. Quelles ressources coûtent même inutilisées ?}\\
Les \textbf{IP publiques statiques}, le \textbf{Load Balancer}, les \textbf{disques managés}, la \textbf{base
SQL provisionnée} et le \textbf{Storage} facturent en continu. Le compute des VM ne coûte que si la VM est
allumée (une VM \textbf{désallouée} ne facture plus de compute, mais son disque OS reste facturé).

\textbf{2. Que peut-on arrêter hors période de formation ?}\\
\textbf{Désallouer les VM} (gain principal) ; éventuellement supprimer le Load Balancer et les IP publiques
(recréables par script). SQL Basic et Storage sont peu coûteux : à laisser ou supprimer via le RG.

\textbf{3. Pourquoi les tags pour analyser les coûts ?}\\
Ils permettent d'\textbf{attribuer} la dépense (application, environnement, équipe), de \textbf{filtrer} dans
Cost Management et de définir des \textbf{budgets ciblés}.

\textbf{4. Réduire vs optimiser les coûts ?}\\
\textbf{Réduire} = baisser la dépense, parfois au détriment du service. \textbf{Optimiser} = meilleur rapport
\textbf{valeur/coût} (bon dimensionnement, suppression du gaspillage) \textbf{sans dégrader la valeur métier}.
La cible FinOps est l'\textbf{optimisation}.

\screenshot{cost-analysis.png}{Cost Management \arr{} Cost analysis sur \code{rg-shopeasy-dev} : peu de données (la facturation ne remonte qu'après 24-48 h).}

\begin{etatbox}
\begin{itemize}
  \item Estimation $\approx$ \textbf{51,5 \$/mois} (24/7) ; postes principaux identifiés (LB, IP, compute).
  \item 3 recommandations FinOps + réponses aux 4 questions.
  \item \textbf{Prêt pour l'Atelier 13 — analyse de disponibilité.}
\end{itemize}
\end{etatbox}


% #####################################################################
%  ATELIER 13
% #####################################################################
\section{Atelier 13 — Analyse de disponibilité}

\begin{fiche}
\textbf{\eTarget~Objectif} — évaluer la robustesse de l'architecture face à des incidents simples.

\textbf{\ePackage~Livrable attendu} — tableau d'analyse d'incidents et proposition d'évolution de l'architecture.
\end{fiche}

\subsection{Scénarios d'incident}
\begin{ltable}{Q[l,0.7] X[l,1.2] Q[l,0.6] X[l,1.3]}
Scénario & Impact & Utilisateurs & Solution technique \\
\textbf{Une VM web tombe en panne} & La sonde du LB la retire ; le service continue sur l'autre VM (capacité ÷ 2). \textbf{Pas d'interruption}. & Aucun en temps normal & \textbf{Déjà couvert} (2 VM + LB + sonde). 3\textsuperscript{e} instance / autoscale pour la capacité, multi-zone pour la résilience \\
\textbf{Le répartiteur est mal configuré} & Sonde/règle/backend erronés \arr{} trafic mal routé ou \textbf{rejeté}, même si les VM sont saines & Tous & Vérifier sonde/règle/backend ; \textbf{IaC} reproductible ; alertes sur la santé du LB ; tester après chaque changement \\
\textbf{La base devient indisponible} & L'app ne peut plus lire/écrire (commandes, factures) \arr{} \textbf{erreurs fonctionnelles} malgré un web up & Tous (données) & SLA \textbf{HA}, \textbf{PITR}, \textbf{failover group} / géo-réplication, monitoring SQL + alertes \\
\textbf{Le Storage est mal configuré} & Documents \textbf{inaccessibles} ou \textbf{exposés} (si public) \arr{} perte d'accès ou fuite & Tous (documents) & Accès privé (déjà), versioning + soft-delete (déjà), \textbf{Private Endpoint} + RBAC, alertes \\
\textbf{Une zone est indisponible} & Les 2 VM étant \textbf{sans zone}, risque de \textbf{perte simultanée} & Tous & \textbf{Déploiement multi-zone} des VM + LB zone-redundant ; base et stockage redondés zone/région \\
\end{ltable}

\subsection{Synthèse — l'architecture suffit-elle pour une application critique ?}
\textbf{Non, pas en l'état.} Points forts déjà acquis :
\begin{itemize}
  \item redondance applicative (2 VM + Load Balancer avec sonde) \arr{} plus de SPOF applicatif ;
  \item \textbf{services managés} pour la base et les documents \arr{} administration et durabilité déléguées ;
  \item segmentation réseau + NSG, supervision de base (métriques + alerte).
\end{itemize}
Limites pour du \textbf{critique} : pas de \textbf{multi-zone} ; base SQL \textbf{Basic} sans HA avancée ;
stockage \textbf{LRS} ; \textbf{aucun test de reprise} formalisé ; IP publiques sur les VM.

\textbf{Évolutions nécessaires (par priorité) :}
\begin{enumerate}
  \item \textbf{Résilience multi-zone} : VM réparties sur $\geq 2$ zones + Load Balancer zone-redundant.
  \item \textbf{Données hautement disponibles} : SQL avec HA (Standard/Premium ou GP) + \textbf{failover group} géo + PITR testé ; stockage \textbf{ZRS/GRS} + Private Endpoint.
  \item \textbf{Sauvegarde \& reprise} : politique formalisée + \textbf{tests de restauration} réguliers (DR drills), \textbf{RPO/RTO} définis.
  \item \textbf{Supervision avancée} : alertes santé LB, DTU SQL, disponibilité HTTP ; Application Insights.
  \item \textbf{Sécurité \& exploitation} : Application Gateway + WAF, Azure Bastion, \textbf{Infrastructure as Code} + pipeline CI/CD.
\end{enumerate}

\begin{etatbox}
\begin{itemize}
  \item 5 scénarios d'incident analysés (impact, utilisateurs, solution).
  \item Verdict + feuille de route d'évolution vers une cible critique.
  \item \textbf{Prêt pour l'Atelier 14 — analyse de sécurité.}
\end{itemize}
\end{etatbox}


% #####################################################################
%  ATELIER 14
% #####################################################################
\section{Atelier 14 — Analyse de sécurité}

\begin{fiche}
\textbf{\eTarget~Objectif} — identifier les risques de sécurité et proposer des mesures correctives adaptées à Azure.

\textbf{\ePackage~Livrable attendu} — matrice de risques et plan d'actions sécurité classé par priorité.
\end{fiche}

\subsection{Matrice de risques}
\begin{tblr}{
  width=\linewidth, colspec={Q[l,0.6] X[l,0.9] Q[c,0.4] X[l,1.2] Q[c,0.2]},
  row{1}={bg=azuredark,fg=white,font=\sffamily\bfseries\scriptsize},
  row{2-Z}={font=\scriptsize}, row{even}={bg=rowalt},
  hline{1,Z}={1.1pt,azuredark}, hline{2}={0.5pt,azuredark!45},
  colsep=4pt, rowsep=3.4pt, stretch=1.15,
}
Risque & Impact & Prob. & Mesure corrective & État \\
\textbf{SSH ouvert à Internet} & Intrusion, force brute, compromission de VM & Élevée & NSG SSH limité à l'IP admin ; cible : Bastion, pas d'IP publique, MFA & \eCheck \\
\textbf{Compte admin partagé} & Perte de traçabilité, audit impossible & Moyenne & Comptes nominatifs Entra ID + RBAC moindre privilège + MFA & \eWarn \\
\textbf{Stockage public} & Fuite de documents clients (RGPD) & Faible & \code{allow-blob-public-access=false} + Private Endpoint + RBAC & \eCheck \\
\textbf{Base exposée} & Vol / altération des données métier & Faible & 0 règle de pare-feu ; Private Endpoint + auth Entra ID + auditing & \eCheck \\
\textbf{Absence de logs/alertes} & Incidents non détectés & Moyenne & Azure Monitor (métriques + alerte) + logs + Defender & \eCheck \\
\textbf{Droits excessifs (Owner)} & Erreur ou compromission à fort impact & Moyenne & RBAC ciblés (Reader/Contributor), revue, PIM & \eWarn \\
\textbf{Chiffrement / TLS faible} & Interception, données en clair & Faible & HTTPS-only + TLS 1.2, chiffrement au repos (TDE/SSE) & \eCheck \\
\end{tblr}
{\small \eCheck~déjà en place dans le TP \qquad \eWarn~à mettre en place (gouvernance).}

\subsection{Plan d'actions sécurité (par priorité)}

\textbf{\eRed~Priorité 1 — Gouvernance des accès (immédiat, faible coût)}
\begin{itemize}
  \item \textbf{Comptes nominatifs} Entra ID (fin du compte partagé) + \textbf{MFA} généralisé.
  \item \textbf{RBAC au moindre privilège} : limiter \code{Owner}, attribuer \code{Reader}/\code{Contributor} ; revue périodique.
  \item Vérifier les NSG : SSH restreint, SQL 1433 depuis le subnet web uniquement.
\end{itemize}

\textbf{\eOrange~Priorité 2 — Réduction de la surface d'exposition (court terme)}
\begin{itemize}
  \item \textbf{Azure Bastion} + suppression des IP publiques des VM.
  \item \textbf{Private Endpoint} pour Azure SQL et Storage.
  \item \textbf{HTTPS obligatoire} + \textbf{Application Gateway / WAF} en frontal.
\end{itemize}

\textbf{\eYellow~Priorité 3 — Détection, résilience et hygiène (moyen terme)}
\begin{itemize}
  \item \textbf{Microsoft Defender for Cloud} + \textbf{NSG Flow Logs} + auditing Azure SQL.
  \item \textbf{Sauvegardes formalisées} + tests de restauration (RPO/RTO).
  \item \textbf{Suppression des ressources inutiles} et \textbf{Infrastructure as Code}.
\end{itemize}

\subsection{Synthèse}
Le modèle est celui de la \textbf{responsabilité partagée} : Azure sécurise l'infrastructure, \textbf{ShopEasy
reste responsable} de la configuration, des identités, des accès et des données. Les bases sont posées (NSG
restrictifs, base et stockage non exposés, TLS/chiffrement, supervision) ; les chantiers prioritaires restants
sont la \textbf{gouvernance des identités} (comptes nominatifs, MFA, moindre privilège) et la \textbf{réduction
de l'exposition} (Bastion, Private Endpoint, WAF).

\begin{etatbox}
\begin{itemize}
  \item Matrice de risques (impact, probabilité, mesure, état réel) et plan d'actions priorisé P1/P2/P3.
  \item \textbf{Prêt pour l'Atelier 15 — note de recommandations DSI.}
\end{itemize}
\end{etatbox}


% #####################################################################
%  ATELIER 15 — NOTE DSI
% #####################################################################
\section{Atelier 15 — Note de recommandations DSI}

\begin{tcolorbox}[enhanced, sharp corners, boxrule=0pt, leftrule=4pt, colback=azurepale, colframe=azuredark,
   left=4.5mm,right=4.5mm,top=3.6mm,bottom=3.6mm, before skip=10pt, after skip=12pt]
\textbf{Destinataire} — Direction des Systèmes d'Information (DSI)\\[2pt]
\textbf{Objet} — proposition d'architecture cloud cible pour l'application de gestion de commandes ShopEasy\\[2pt]
\textbf{Auteurs (binôme)} — Louis \textsc{Scarfone} \& Maxence \textsc{Bourragué}
\end{tcolorbox}

\subsection{Contexte et enjeux}
ShopEasy exploite une application interne de gestion de commandes hébergée sur \textbf{un serveur physique
unique} (Apache/Nginx + PHP/Node.js + MySQL colocalisés, documents en local, sauvegardes manuelles, aucune
supervision, compte administrateur partagé). Cette situation présente des \textbf{risques majeurs} : point de
défaillance unique, sécurité faible, exploitation artisanale, coûts non maîtrisés. La direction souhaite une
première proposition Azure améliorant la \textbf{disponibilité}, \textbf{séparant les couches}, \textbf{réduisant
les risques d'administration} et offrant une \textbf{visibilité financière}.

\subsection{Architecture cible proposée}
\screenshot{architecture.png}{Architecture cible : Internet \arr{} Load Balancer \arr{} 2 VM web (\code{snet-web}) ; VM \arr{} SQL Database et Storage en accès privé (\code{snet-data}) ; gouvernance Entra ID/RBAC, Azure Monitor, Cost Management.}
L'architecture \textbf{sépare les couches} (réseau, applicatif, données, stockage), supprime le \textbf{point de
défaillance unique applicatif} (2 VM + répartiteur de charge) et \textbf{n'expose à Internet que le Load
Balancer}. La base et les documents sont en \textbf{accès privé}.

\subsection{Services Azure retenus et leur rôle}
\begin{ltable}{Q[l,0.7] X[l,1.2] Q[c,0.4]}
Service & Rôle & Modèle \\
Resource Group & Cycle de vie commun, droits, coûts & Gouvernance \\
Virtual Network + subnets & Réseau privé isolé, segmentation web/data/admin & IaaS \\
Network Security Groups & Filtrage des flux (80/443, 22, 1433) & IaaS \\
2 × Virtual Machines (\code{B2ats\_v2}) & Couche applicative redondante (Nginx) & IaaS \\
Azure Load Balancer & Répartition de charge + sonde de santé & IaaS \\
Azure SQL Database (Basic) & Base relationnelle managée, non exposée & PaaS \\
Storage Account (Blob) & Documents clients durables, privés, versionnés & PaaS \\
Microsoft Entra ID + RBAC & Identités et droits (moindre privilège) & Gouvernance \\
Azure Monitor & Métriques, logs, alertes & Gouvernance \\
Cost Management + tags & Suivi et pilotage des coûts & Gouvernance \\
\end{ltable}

\subsection{Gains attendus}
\begin{itemize}
  \item \textbf{Disponibilité} : redondance applicative + LB \arr{} plus de panne bloquante sur une VM.
  \item \textbf{Sécurité} : couches cloisonnées, base et stockage non exposés, SSH restreint, chiffrement/TLS.
  \item \textbf{Exploitation} : services managés (base, stockage), supervision centralisée et alertes.
  \item \textbf{Données} : durabilité, versioning, soft-delete, sauvegardes intégrées de la base.
  \item \textbf{Coûts} : passage en \textbf{OPEX}, visibilité par tags, leviers d'optimisation.
\end{itemize}

\subsection{Sécurité — mesures et risques résiduels}
\textbf{En place :} NSG restrictifs (SSH limité, SQL interne au subnet web), base \textbf{sans accès public}
(0 règle de pare-feu), \textbf{stockage public bloqué}, \textbf{TLS 1.2} + chiffrement au repos, supervision et
alerte. \textbf{Risques résiduels :} compte d'administration à rendre nominatif (+ MFA, RBAC) ; IP publiques
des VM à supprimer (\arr{} Bastion) ; passage en \textbf{Private Endpoint} pour SQL et Storage ; HTTPS/WAF via
Application Gateway.

\subsection{Disponibilité — niveau obtenu et améliorations}
\textbf{Obtenu :} tolérance à la panne d'\textbf{une} VM, base et stockage managés et redondés localement.
\textbf{Limites \& améliorations :} déploiement \textbf{mono-zone} (contrainte Azure for Students) \arr{} cible
\textbf{multi-zone} ; base \textbf{Basic} \arr{} niveau supérieur + failover ; stockage \textbf{LRS} \arr{}
ZRS/GRS ; \textbf{tests de restauration} et objectifs RPO/RTO à formaliser.

\subsection{Coûts (FinOps)}
Estimation \textbf{vérifiée via l'API Azure Retail Prices} (juin 2026, pay-as-you-go, 24/7) :
\begin{ltable}{X[l,1] Q[c,0.4]}
Poste & Coût (\$/mois) \\
2 × VM \code{B2ats\_v2} & 14,20 \\
Load Balancer Standard & 18,25 \\
3 × IP publique Standard & 10,95 \\
Base SQL Basic & 4,90 \\
2 × disque HDD S4 & 3,07 \\
Storage + Monitoring & $\approx$ 0,12 \\
\textbf{Total} & \textbf{$\approx$ 51,5} \\
\end{ltable}
\textbf{Recommandations FinOps :} désallouer les VM hors usage ; supprimer les ressources qui facturent à vide
(LB, IP) hors période d'activité ; réduire les IP publiques (3 \arr{} 1) ; dimensionner les niveaux ; piloter
via tags + budgets/alertes.

\subsection{Plan d'action priorisé}
\begin{ltable}{Q[l,0.4] X[l,1.4]}
Échéance & Actions \\
\textbf{Court terme} & Comptes nominatifs + MFA, RBAC moindre privilège ; budgets/alertes de coût ; HTTPS ; désallocation/nettoyage des ressources de test. \\
\textbf{Moyen terme} & Azure Bastion + suppression des IP publiques ; Private Endpoint SQL/Storage ; Application Gateway + WAF ; multi-zone ; SQL HA + failover ; sauvegardes testées ; supervision avancée. \\
\textbf{Long terme} & Infrastructure as Code (Terraform/Bicep) + pipeline CI/CD ; éventuelle bascule vers App Service (PaaS). \\
\end{ltable}

\subsection{Limites de la proposition}
Architecture \textbf{volontairement progressive et pédagogique} : mono-zone, niveaux minimaux (Basic/LRS), VM
avec IP publiques pour faciliter le TP, accès aux données par pare-feu plutôt que Private Endpoint. Elle
\textbf{démontre les briques fondamentales} mais \textbf{n'est pas une cible de production critique} : les
évolutions du plan d'action sont nécessaires avant une mise en production.

\begin{notebox}[Contraintes d'environnement (Azure for Students)]
Régions restreintes par policy (déploiement en \textbf{Sweden Central}), tailles de VM limitées
(\textbf{\code{Standard\_B2ats\_v2}} au lieu de \code{B1s}), déploiement zonal indisponible, providers à
enregistrer manuellement — éléments à anticiper pour une souscription d'entreprise.
\end{notebox}

\subsection*{Conclusion}
La migration proposée transforme un serveur unique fragile en une architecture \textbf{segmentée, redondée,
supervisée et au coût maîtrisé}, en s'appuyant sur des services Azure adaptés à chaque besoin. Elle répond aux
objectifs de la DSI tout en traçant une \textbf{feuille de route claire} vers une cible de production sécurisée
et hautement disponible.


% #####################################################################
%  QUIZ
% #####################################################################
\section{Quiz de validation}

\begin{enumerate}[leftmargin=2.2em]
\item \textbf{Quel service Azure permet de créer un réseau privé logique ?}\\ Azure Virtual Network (VNet).
\item \textbf{Quel service Azure permet de créer une machine virtuelle ?}\\ Azure Virtual Machines.
\item \textbf{Différence entre une région et une zone de disponibilité ?}\\ Une \textbf{région} est une zone géographique ; une \textbf{zone de disponibilité} est un ensemble de datacenters physiquement séparés au sein d'une région.
\item \textbf{À quoi sert un Resource Group ?}\\ À regrouper des ressources partageant un cycle de vie, pour les administrer, taguer, suivre en coût et supprimer ensemble.
\item \textbf{Pourquoi utiliser plusieurs subnets ?}\\ Pour séparer les rôles (web/data/admin) et appliquer des règles de sécurité différentes par couche.
\item \textbf{À quoi sert un Network Security Group ?}\\ À filtrer le trafic réseau entrant et sortant (port, protocole, source/destination).
\item \textbf{Pourquoi ne pas ouvrir SSH à tout Internet ?}\\ Pour réduire la surface d'attaque (scans et force brute) ; on le limite à l'IP admin ou via Azure Bastion.
\item \textbf{Quel service remplace le stockage local de documents ?}\\ Azure Storage Account (Blob Storage).
\item \textbf{Quel service remplace une base SQL installée sur un serveur ?}\\ Azure SQL Database (base managée / PaaS).
\item \textbf{Pourquoi utiliser un répartiteur de charge ?}\\ Pour distribuer le trafic et réduire l'impact de la panne d'un serveur (suppression du SPOF applicatif).
\item \textbf{Différence entre Azure Load Balancer et Application Gateway ?}\\ Load Balancer = couche 4 (TCP/UDP) ; Application Gateway = couche 7 (HTTP/HTTPS) avec WAF, terminaison TLS, routage par chemin.
\item \textbf{Que signifie PaaS ?}\\ Platform as a Service (plateforme managée : Azure gère l'OS/runtime, le client gère le code et les données).
\item \textbf{Azure SQL Database est-il IaaS ou PaaS ?}\\ PaaS.
\item \textbf{Pourquoi utiliser Azure Monitor ?}\\ Pour collecter métriques et logs, déclencher des alertes et obtenir une visibilité opérationnelle.
\item \textbf{Que surveiller sur une VM web ?}\\ CPU, mémoire, disque, disponibilité HTTP, erreurs, activité réseau.
\item \textbf{Que permet Azure Cost Management ?}\\ Suivre, analyser et optimiser les coûts (budgets, alertes, répartition par tags).
\item \textbf{Pourquoi taguer les ressources ?}\\ Pour identifier propriétaire/projet/environnement/centre de coût et filtrer, attribuer et analyser les coûts.
\item \textbf{Quel risque pose un compte administrateur partagé ?}\\ Perte de traçabilité et droits trop larges (audit et conformité compromis).
\item \textbf{Pourquoi éviter un Storage Account public par défaut ?}\\ Risque d'exposition de données sensibles (fuite RGPD).
\item \textbf{Quelle mesure réduit l'impact d'une panne de VM ?}\\ Plusieurs VM derrière un répartiteur de charge (avec sonde de santé).
\item \textbf{Pourquoi séparer la couche web et la couche données ?}\\ Pour réduire les mouvements latéraux en cas de compromission et appliquer la segmentation.
\item \textbf{Quelle information doit apparaître dans une note DSI ?}\\ Architecture, coûts, risques, gains et plan d'action.
\item \textbf{Citez deux optimisations de coût.}\\ Arrêt/désallocation des VM inutilisées et choix de tailles adaptées (+ tags, budgets, réservations).
\item \textbf{Citez deux mesures de sécurité prioritaires.}\\ RBAC + moindre privilège (avec MFA) et NSG restrictifs (+ chiffrement, suppression des accès publics).
\item \textbf{Quelle évolution permettrait d'automatiser le déploiement ?}\\ L'Infrastructure as Code (Terraform ou Bicep), idéalement avec un pipeline CI/CD.
\end{enumerate}

\vspace{6pt}
\begin{tcolorbox}[enhanced, sharp corners, boxrule=0pt, leftrule=4pt, colback=okgreenbg, colframe=okgreen,
   left=4.5mm,right=4.5mm,top=3.6mm,bottom=3.6mm]
{\sffamily\bfseries\color{okgreen}\eParty~Fin du dossier} — les 15 ateliers et le quiz de validation sont
complets : analyse de l'existant, choix des services, architecture cible, déploiement Azure CLI (réseau,
VM, load balancer, stockage, base SQL, supervision), FinOps, disponibilité, sécurité et note DSI.
\end{tcolorbox}

\end{document}
